\chapter{Line bundles, metrics, and model measures}
\index{metric!model}\index{Chambert--Loir measure}
\label{chap:metrics}

The preceding chapter turned vertical divisors into functions. A vertical
line bundle does more: it changes the norm of every local section of a line
bundle on the generic fiber. Relative nefness then becomes semipositivity,
and intersection numbers on the special fiber become a measure on the
Berkovich space.

This chapter constructs that model core. It is the common domain on which the
global BFJ theory, the differential CLD theory, and later variants must agree.
We begin with continuous metrics, pass to algebraic and formal models,
introduce Fubini--Study metrics, and then define semipositivity. The final
sections construct Chambert--Loir measures and prove their elementary
intersection-theoretic properties.

It is worth saying at the outset what is being replaced. On a complex
manifold, a smooth Hermitian metric on a line bundle has a curvature form,
and one may differentiate a metric weight twice to obtain a measure. Nothing
of the kind is available yet on a Berkovich space: no differential calculus
has been constructed at this point of the book, and the reader should resist
the temptation to imagine one. What is available instead is the inverse
system of models built in \cref{chap:models} and the vertical divisors and
skeleta of \cref{chap:skeleta}. A metric will therefore be produced by an
algebro-geometric datum, its positivity will be tested by intersection
numbers against curves in a special fiber, and its ``curvature'' will be an
atomic measure whose masses are intersection numbers. Only in
\cref{part:cld} does a genuine differential calculus appear, and one of the
main comparison theorems of the book, in \cref{chap:compare-ma}, states that
it recovers exactly the measures constructed here.

Four descriptions of the same class of metrics will occur, and it is useful
to have their names fixed before they are constructed:
\begin{enumerate}
\item a metric induced by a line bundle \(\mathcal L\) on an algebraic model
\(\mathcal X\), possibly after extracting a root: an \emph{algebraic model
metric};
\item a metric induced by a line bundle on an admissible formal model: a
\emph{formal metric};
\item a metric admitting local frames of norm one on an analytic-domain
cover: a \emph{piecewise linear metric};
\item a metric assembled from finitely many global sections and real
weights: a \emph{Fubini--Study metric}.
\end{enumerate}
The first three agree on a proper scheme, by a theorem of Gubler--Martin
recalled and proved in \cref{sec:metrics-formal-pl}; the fourth is a special
case of the first when the weights lie in the divisible value group, and is
otherwise a uniform limit of such. Keeping the four names distinct is not
pedantry. Over a trivially valued field the equivalence fails in the form
stated here, and \cref{part:bj} builds its theory on the fourth description
alone.

\subsection*{Plan and dependencies}

\Cref{sec:metrics-continuous} constructs continuous metrics, additive
weights, and the correspondence between metrics and functions once a
reference metric is fixed; the change-of-reference bookkeeping is done there
once and for all, with the signs displayed.
\Cref{sec:metrics-model} constructs algebraic model metrics, proves that
they are continuous, and computes the difference of two of them as the model
function of a vertical divisor; it also records the dictionary between
vertical divisors and model functions with all normalizations explicit.
\Cref{sec:metrics-formal-pl} compares algebraic, formal, and piecewise
linear metrics over an arbitrary complete nontrivially valued field.
\Cref{sec:metrics-fs} treats Fubini--Study metrics and computes several of
them explicitly.
\Cref{sec:metrics-model-functions} isolates the space of model functions,
proves that it is a lattice separating points, and deduces its uniform
density in the continuous functions; this density theorem, due to Gubler, is
the analytic engine of every subsequent approximation argument in the book.
\Cref{sec:metrics-semipositive} defines semipositivity through nef models
and proves its stability under pullback, tensor product, maximum, and
uniform limits.
\Cref{sec:metrics-positivity-interface} separates the three positivity
languages that meet here and states precisely what is, and what is not,
asserted about their comparison.
\Cref{sec:metrics-cl-measure} constructs Chambert--Loir measures at model
level, verifies the total mass \((L_1\cdots L_n)\), and computes four
examples in full.
\Cref{sec:metrics-continuous-cl} extends the construction to continuous
semipositive metrics, and \cref{sec:metrics-operations} records the
functorial calculus.

The prerequisites are \cref{chap:models} for models, vertical divisors, and
intersection numbers, \cref{chap:berkovich} for analytification, centers,
and reduction, and \cref{chap:skeleta} for skeleta and model functions. From
\cref{chap:model-calculations} we use only the two worked calculi---graph
Laplacians and toric Alexandrov measures---and we use them to check answers,
never to obtain them. Readers following the complex/Kähler route may read
\crefrange{sec:metrics-continuous}{sec:metrics-fs} and then pass to
\cref{part:complex}.

\subsection*{Standing hypotheses}

Unless explicitly stated otherwise, \(K\) is a complete nontrivially valued
field and \(X\) is a proper \(K\)-scheme. In the sections involving
intersection numbers on a special fiber, \(K\) is discretely valued,
\(R=K^\circ\), a uniformizer \(\pi\) is normalized by \(v(\pi)=1\), and
\(\widetilde K\) is the residue field. The trivially valued Fubini--Study
theory is developed separately in \cref{part:bj}.

The reason for the two regimes is not fastidiousness. Over a general
complete nontrivially valued field the valuation ring need not be
noetherian, so Chow-theoretic intersection numbers on models are not
available in the form used in \cref{chap:models}; conversely, formal
geometry is available over any such field and is the natural setting for the
comparison theorems of \cref{sec:metrics-formal-pl}. Every statement below
therefore carries its field hypothesis visibly. A statement proved for a
discretely valued field is never silently used over an arbitrary one.

Finally, a word on the word \emph{semipositive}. It will be defined by
relative nefness of a model line bundle, then extended to continuous metrics
by uniform approximation. This is Zhang's definition. It is a global,
approximation-theoretic notion. It is not, at this stage of the book, a
local differential condition, and no local characterization of it is used
before \cref{part:cld} supplies one with its own hypotheses.

\section{Continuous metrics and additive weights}
\label{sec:metrics-continuous}

Let \(L\) be a line bundle on \(X\), and denote its analytification by
\(L^{\mathrm{an}}\). Concretely, \(L^{\mathrm{an}}\) is the analytification
of the total space of \(L\), but for our purposes only its fibers matter.
For an analytic point \(x\in X^{\mathrm{an}}\) with support
\(\rho_X(x)\in X\) and completed residue field \(\mathcal H(x)\), the fiber
of \(L\) at \(x\) is the one-dimensional \(\mathcal H(x)\)-vector space
\[
 L(x)=L_{\rho_X(x)}\otimes_{\mathcal O_{X,\rho_X(x)}}\mathcal H(x).
\]
A local section \(s\) of \(L\) near \(\rho_X(x)\) has a well-defined value
\(s(x)\in L(x)\). Two facts from \cref{chap:berkovich} are used constantly
and without further comment: \(\mathcal H(x)\) carries an absolute value
extending that of \(K\), and \(x\mapsto|f(x)|\) is continuous for every
analytic function \(f\). A metric is a rule assigning a norm to each of the
lines \(L(x)\), subject to a continuity requirement.

\begin{definition}[Continuous metric]
A \emph{continuous metric} \(\|\cdot\|\) on \(L^{\mathrm{an}}\) assigns to
each analytic point \(x\) a norm on the one-dimensional
\(\mathcal H(x)\)-vector space \(L(x)\), continuously in local
trivializations. Thus, for a local frame \(\tau\), the function
\(x\mapsto\|\tau(x)\|\) is positive and continuous, and
\[
                    \|f\tau(x)\|=|f(x)|\,\|\tau(x)\|
\]
for every analytic function \(f\).
\end{definition}

Two remarks on the definition. First, the continuity requirement is
local and frame-dependent in appearance only: if \(\tau'=g\tau\) is another
frame on a smaller domain, with \(g\) an invertible analytic function there,
then \(\|\tau'\|=|g|\,\|\tau\|\), and \(|g|\) is continuous and nowhere
zero. Continuity in one frame is therefore continuity in all frames.
Second, the homogeneity condition \(\|f\tau(x)\|=|f(x)|\,\|\tau(x)\|\) is
exactly the statement that \(\|\cdot\|\) restricts to a norm on each
\(\mathcal H(x)\)-line \(L(x)\), so no positivity or triangle inequality has
to be imposed separately: a norm on a one-dimensional vector space over a
valued field is determined by its value on one nonzero vector, and any
positive value is allowed.

\subsection{Additive weights and the cocycle rule}

It is often clearer to use additive weights. The weight of a metric in the
frame \(\tau\) is
\[
                         \phi_\tau=-\log\|\tau\|.
\]
If \(\tau'=g\tau\), then
\[
                         \phi_{\tau'}=\phi_\tau-\log|g|.
\]
This transformation rule is exactly what is needed for the local weights to
define a metric. The sign is the one fixed in \cref{sec:appJ-metrics}: a
\emph{larger} weight means a \emph{smaller} norm, in accordance with
\eqref{eq:appJ-na-weight}. We record the elementary but foundational
statement that weights and metrics are interchangeable data.

\begin{lemma}
\label{lem:metrics-weight-cocycle}
Let \(\{V_i\}_{i\in I}\) be a covering of \(X^{\mathrm{an}}\) by open
subsets, and for each \(i\) let \(\tau_i\) be a frame of \(L^{\mathrm{an}}\)
on \(V_i\), with transition functions \(g_{ij}\) determined by
\(\tau_j=g_{ij}\tau_i\) on \(V_i\cap V_j\). Suppose given continuous
functions \(\phi_i:V_i\to\mathbf R\) satisfying
\begin{equation}
\label{eq:metrics-weight-cocycle}
 \phi_j=\phi_i-\log|g_{ij}|
 \qquad\text{on }V_i\cap V_j .
\end{equation}
Then there is a unique continuous metric \(\|\cdot\|\) on
\(L^{\mathrm{an}}\) with \(-\log\|\tau_i\|=\phi_i\) for every \(i\).
Conversely, every continuous metric arises in this way from its own local
weights.
\end{lemma}

\begin{proof}
Uniqueness is clear, since a metric is determined by the norm of one
nonvanishing local section at each point. For existence, let
\(x\in X^{\mathrm{an}}\) and let \(\xi\in L(x)\). Choose \(i\) with
\(x\in V_i\) and write \(\xi=c\,\tau_i(x)\) with \(c\in\mathcal H(x)\); this
is possible and \(c\) is unique because \(\tau_i(x)\) is a basis of the
one-dimensional \(\mathcal H(x)\)-vector space \(L(x)\). Set
\[
 \|\xi\|:=|c|\,\rme^{-\phi_i(x)} .
\]
If also \(x\in V_j\), then \(\tau_j(x)=g_{ij}(x)\tau_i(x)\), so the
coefficient of \(\xi\) in the frame \(\tau_j\) is \(c\,g_{ij}(x)^{-1}\), and
the corresponding value is
\[
 |c|\,|g_{ij}(x)|^{-1}\rme^{-\phi_j(x)}
 =|c|\,|g_{ij}(x)|^{-1}\rme^{-\phi_i(x)}\,|g_{ij}(x)|
 =|c|\,\rme^{-\phi_i(x)},
\]
where the middle equality is \eqref{eq:metrics-weight-cocycle} in the form
\(\rme^{-\phi_j}=\rme^{-\phi_i}|g_{ij}|\). Hence \(\|\xi\|\) is well
defined. Homogeneity holds because \(|fc|=|f|\,|c|\) in \(\mathcal H(x)\),
and continuity holds in the frame \(\tau_i\) by construction, hence in every
frame by the remark preceding this lemma. The converse statement is the
definition of \(\phi_{\tau}\) together with the transformation rule.
\end{proof}

The cocycle rule \eqref{eq:metrics-weight-cocycle} is the additive shadow of
the multiplicative cocycle \(\{g_{ij}\}\) that defines \(L\). It is worth
noticing that no continuity of \(g_{ij}\) beyond that of \(|g_{ij}|\) is
used; this is why the same formalism will accept model frames, formal
frames, and global sections without modification.

\subsection{Relative weights: metrics as a torsor}

Fix a reference metric \(\|\cdot\|_0\). For any other metric
\(\|\cdot\|\), the relative weight
\[
 \phi=-\log\frac{\|\tau\|}{\|\tau\|_0}
\]
is independent of the frame and hence is a continuous function on
\(X^{\mathrm{an}}\). Conversely,
\[
                        \|\cdot\|=\|\cdot\|_0\rme^{-\phi}
\]
defines a metric for every continuous \(\phi\). Thus the space of continuous
metrics on a fixed line bundle is an affine space modeled on
\(C^0(X^{\mathrm{an}})\). We state this precisely, because the whole book
depends on passing back and forth between metrics and functions, and because
the change-of-reference rule has to be memorized with its sign.

\begin{proposition}
\label{prop:metrics-torsor}
Let \(L\) be a line bundle on the proper \(K\)-scheme \(X\), and let
\(\mathcal M(L)\) denote the set of continuous metrics on
\(L^{\mathrm{an}}\), assumed nonempty. Then the following hold.
\begin{enumerate}
\item The formula
\((\phi,\|\cdot\|_0)\mapsto\|\cdot\|_0\rme^{-\phi}\) defines a simply
transitive action of the additive group \(C^0(X^{\mathrm{an}})\) on
\(\mathcal M(L)\). Consequently, once a reference \(\|\cdot\|_0\) is fixed,
\[
 C^0(X^{\mathrm{an}})\xrightarrow{\ \sim\ }\mathcal M(L),
 \qquad
 \phi\longmapsto\|\cdot\|_0\rme^{-\phi},
\]
is a bijection, and we write \(\phi=\|\cdot\|/\|\cdot\|_0\) additively.
\item If the reference is changed from
\(\|\cdot\|_0\) to \(\|\cdot\|_0'=\|\cdot\|_0\rme^{-\chi}\), with
\(\chi\in C^0(X^{\mathrm{an}})\), then the relative weight of a fixed metric
changes by
\[
 \phi'=\phi-\chi .
\]
In particular differences of relative weights, and hence all statements of
the form \(\phi_1\le\phi_2\), \(\phi_1=\phi_2+c\), or
\(\sup(\phi_1-\phi_2)\le\varepsilon\), are independent of the reference.
\item If \(\tau\) is a frame of \(L\) on an
analytic domain \(V\), then on \(V\)
\[
 \phi=\phi_\tau-\phi_{\tau,0},
\]
the difference of the two absolute weights in that frame. The right-hand
side is unchanged if \(\tau\) is replaced by \(g\tau\).
\end{enumerate}
\end{proposition}

\begin{proof}
(1) Let \(\|\cdot\|_0\in\mathcal M(L)\) and
\(\phi\in C^0(X^{\mathrm{an}})\). The assignment
\(\xi\mapsto\rme^{-\phi(x)}\|\xi\|_0\) is again a norm on each line
\(L(x)\), since \(\rme^{-\phi(x)}>0\), and it is continuous in every frame
because \(\phi\) is continuous. The action is additive in \(\phi\), and
\(\phi=0\) acts trivially. Given two metrics \(\|\cdot\|_1,\|\cdot\|_2\),
the function
\(\phi:=-\log(\|\tau\|_1/\|\tau\|_2)\), computed in any local frame
\(\tau\), is independent of \(\tau\) by (3)
below and is continuous; it is the unique element of
\(C^0(X^{\mathrm{an}})\) carrying \(\|\cdot\|_2\) to \(\|\cdot\|_1\). Hence
the action is simply transitive.

(2) By definition
\(\|\cdot\|=\|\cdot\|_0\rme^{-\phi}\) and
\(\|\cdot\|_0=\|\cdot\|_0'\rme^{\chi}\), so
\(\|\cdot\|=\|\cdot\|_0'\rme^{-(\phi-\chi)}\), which is the assertion.
Any expression involving only differences of relative weights taken with
respect to the same reference is therefore unchanged, since the two copies
of \(\chi\) cancel.

(3) In the frame \(\tau\) one has
\(\|\tau\|=\rme^{-\phi_\tau}\) and
\(\|\tau\|_0=\rme^{-\phi_{\tau,0}}\), whence
\[
 -\log\bigl(\|\tau\|/\|\tau\|_0\bigr)=\phi_\tau-\phi_{\tau,0}.
\]
Replacing
\(\tau\) by \(g\tau\) adds \(-\log|g|\) to both absolute weights, and the
difference is unchanged.
\end{proof}

The reader should notice what \cref{prop:metrics-torsor} does \emph{not}
say. It does not say that there is a canonical continuous metric on \(L\);
there is none, and the choice of a reference is a genuine choice, usually
made by choosing a model. It does not say that a distinguished function is
attached to a metric; only differences of metrics are functions. Every later
notion---model function, psh function, energy---is defined relative to a
reference and transforms by the rule in (2).

\begin{definition}[Uniform distance]
For two continuous metrics on \(L\), define
\[
 d_\infty(\|\cdot\|_1,\|\cdot\|_2)
 =\sup_{X^{\mathrm{an}}}
    \left|\log\frac{\|\tau\|_1}{\|\tau\|_2}\right|.
\]
This is independent of \(\tau\). It is finite because \(X^{\mathrm{an}}\)
is compact.
\end{definition}

Under the bijection of \cref{prop:metrics-torsor}, the uniform distance is
the sup-norm distance of relative weights:
\begin{equation}
\label{eq:metrics-dinfty-sup}
 d_\infty\bigl(\|\cdot\|_0\rme^{-\phi_1},\|\cdot\|_0\rme^{-\phi_2}\bigr)
 =\|\phi_1-\phi_2\|_\infty .
\end{equation}
In particular \(d_\infty\) does not depend on the reference metric, in
accordance with \cref{prop:metrics-torsor}\,(2). The next lemma is used
whenever a metric is produced as a uniform limit, which in this book means
essentially always.

\begin{lemma}
\label{lem:metrics-dinfty-complete}
Let \(X\) be a proper \(K\)-scheme and \(L\) a line bundle carrying at least
one continuous metric. Then \((\mathcal M(L),d_\infty)\) is a complete
metric space, and a \(d_\infty\)-Cauchy sequence of continuous metrics
converges uniformly on \(X^{\mathrm{an}}\), in every local frame, to a
continuous metric.
\end{lemma}

\begin{proof}
By \Cref{thm:berk-proper-compact} the space \(X^{\mathrm{an}}\) is compact
and Hausdorff, so \(C^0(X^{\mathrm{an}})\) with the sup-norm is a Banach
space. Fix a reference \(\|\cdot\|_0\). By \Cref{prop:metrics-torsor}\,(1)
the map \(\phi\mapsto\|\cdot\|_0\rme^{-\phi}\) is a bijection from
\(C^0(X^{\mathrm{an}})\) to \(\mathcal M(L)\), and by
\eqref{eq:metrics-dinfty-sup} it is an isometry from the sup-norm distance
to \(d_\infty\). An isometric image of a complete metric space is complete.
For the last assertion, if \(\phi_j\to\phi\) uniformly then, in a frame
\(\tau\) over an analytic domain \(V\),
\(\|\tau\|_j=\rme^{-\phi_j}\|\tau\|_0\to\rme^{-\phi}\|\tau\|_0\) uniformly
on \(V\), because \(\|\tau\|_0\) is bounded on the compact set
\(\overline V\subset X^{\mathrm{an}}\) and \(t\mapsto\rme^{-t}\) is Lipschitz
on bounded sets.
\end{proof}

\subsection{Tensor product, duality, and pullback}

Tensor product and duality add and negate local weights:
\[
 \phi_{L\otimes M}=\phi_L+\phi_M,\qquad
 \phi_{L^\vee}=-\phi_L.
\]
Pullback composes weights with the analytic map. We spell this out, because
the statements are used in every subsequent chapter and because the
compatibility of the operations with relative weights, rather than with
absolute ones, is what makes them well defined.

\begin{proposition}
\label{prop:metrics-operations}
Let \(L,M\) be line bundles on the proper \(K\)-scheme \(X\), equipped with
continuous metrics \(\|\cdot\|_L,\|\cdot\|_M\), and let \(f:Y\to X\) be a
morphism of proper \(K\)-schemes.
\begin{enumerate}
\item There is a unique continuous metric \(\|\cdot\|_{L\otimes M}\) on
\((L\otimes M)^{\mathrm{an}}\) with
\(\|s\otimes t\|_{L\otimes M}=\|s\|_L\|t\|_M\) for local sections \(s,t\);
in local frames its weight is \(\phi_{L,\tau}+\phi_{M,\sigma}\), and for
relative weights \(\phi_{L\otimes M}=\phi_L+\phi_M\) once the references are
tensored.
\item There is a unique continuous metric on \((L^\vee)^{\mathrm{an}}\)
making the evaluation pairing \(L\otimes L^\vee\to\mathcal O_X\) an isometry
onto the trivial metric; its weight is \(-\phi_L\).
\item There is a unique continuous metric \(f^*\|\cdot\|_L\) on
\((f^*L)^{\mathrm{an}}\) with \(\|f^*s\|=\|s\|_L\circ f^{\mathrm{an}}\), and
its relative weight is \(\phi_L\circ f^{\mathrm{an}}\).
\item For two continuous metrics on the same \(L\), the pointwise minimum of
the norms is a continuous metric, whose relative weight is the pointwise
maximum \(\max\{\phi_1,\phi_2\}\) of the relative weights.
\end{enumerate}
Moreover \(d_\infty\) is subadditive for tensor products and nonincreasing
under pullback:
\[
\begin{gathered}
 d_\infty(\|\cdot\|_1\otimes\|\cdot\|_1',\|\cdot\|_2\otimes\|\cdot\|_2')
 \le d_\infty(\|\cdot\|_1,\|\cdot\|_2)
   +d_\infty(\|\cdot\|_1',\|\cdot\|_2'),\\
 d_\infty(f^*\|\cdot\|_1,f^*\|\cdot\|_2)\le
 d_\infty(\|\cdot\|_1,\|\cdot\|_2).
\end{gathered}
\]
\end{proposition}

\begin{proof}
(1) On the fiber at \(x\) one has
\((L\otimes M)(x)=L(x)\otimes_{\mathcal H(x)}M(x)\), a one-dimensional
\(\mathcal H(x)\)-vector space, and the tensor product of two norms on
one-dimensional spaces is the unique norm with the stated multiplicativity.
If \(\tau,\sigma\) are frames of \(L,M\) over \(V\), then
\(\tau\otimes\sigma\) is a frame of \(L\otimes M\) over \(V\) with
\(\|\tau\otimes\sigma\|=\rme^{-(\phi_{L,\tau}+\phi_{M,\sigma})}\), which is
continuous and positive. The transition rule
\eqref{eq:metrics-weight-cocycle} for \(L\otimes M\) is the sum of those for
\(L\) and \(M\), so \Cref{lem:metrics-weight-cocycle} applies.

(2) The pairing identifies \(L^\vee(x)\) with the \(\mathcal H(x)\)-dual of
\(L(x)\); the dual norm on the dual of a one-dimensional normed space
satisfies \(\|\lambda\|\,\|\xi\|=|\lambda(\xi)|\), whence
\(\|\tau^\vee\|=\|\tau\|^{-1}\) for the dual frame and
\(\phi_{L^\vee,\tau^\vee}=-\phi_{L,\tau}\).

(3) For \(y\in Y^{\mathrm{an}}\) there is a canonical isometric embedding
\(\mathcal H(f^{\mathrm{an}}(y))\hookrightarrow\mathcal H(y)\), by the
definition of the analytification of a morphism in
\cref{sec:berk-functoriality}, and
\((f^*L)(y)=L(f^{\mathrm{an}}(y))\otimes\mathcal H(y)\). Transporting the
norm along this embedding gives the asserted metric; in a frame
\(f^*\tau\), its weight is \(\phi_{L,\tau}\circ f^{\mathrm{an}}\), a
continuous function since \(f^{\mathrm{an}}\) is continuous.

(4) In a frame \(\tau\) the two norms are \(\rme^{-\phi_{1,\tau}}\) and
\(\rme^{-\phi_{2,\tau}}\); their minimum is
\(\rme^{-\max\{\phi_{1,\tau},\phi_{2,\tau}\}}\), and the maximum of two
continuous functions is continuous. The cocycle rule is preserved because
both weights transform by the same additive term \(-\log|g|\), which
commutes with taking maxima:
\(\max\{a-c,b-c\}=\max\{a,b\}-c\). Passing to relative weights with respect
to a common reference gives the stated formula.

The two inequalities follow from \eqref{eq:metrics-dinfty-sup}: for the
first, \(\|(\phi_1+\phi_1')-(\phi_2+\phi_2')\|_\infty\le
\|\phi_1-\phi_2\|_\infty+\|\phi_1'-\phi_2'\|_\infty\); for the second, the
sup of \(|\phi_1-\phi_2|\circ f^{\mathrm{an}}\) over \(Y^{\mathrm{an}}\) is
at most the sup of \(|\phi_1-\phi_2|\) over \(X^{\mathrm{an}}\).
\end{proof}

\begin{example}[Metrics on the trivial bundle are functions]
\label{ex:metrics-trivial-bundle}
Take \(L=\mathcal O_X\) with the reference metric \(\|1\|_0\equiv1\). By
\Cref{prop:metrics-torsor} every continuous metric on
\(\mathcal O_X^{\mathrm{an}}\) is \(\|1\|=\rme^{-\phi}\) for a unique
\(\phi\in C^0(X^{\mathrm{an}})\), and conversely. This identification is
used so often that we shall not mention it: a continuous function on
\(X^{\mathrm{an}}\) \emph{is} a continuous metric on the trivial bundle, and
a model function will turn out to be a model metric on the trivial bundle.
The sign test to remember is that a large \(\phi\) makes the constant
section \(1\) short.

For a concrete instance, let \(f\in K(X)^\times\) be a nonzero rational
function, regarded as a rational section of \(\mathcal O_X\). Then
\(\|f\|=|f|\rme^{-\phi}\), so the metric with weight \(\phi=-\log|f|=v(f)\)
makes \(\|f\|\equiv1\). This is the germ of the model construction of
\cref{sec:metrics-model}: a divisor is a way of prescribing which section
has norm one.
\end{example}

\begin{example}[The standard metric on \(\mathcal O(1)\), in two frames]
\label{ex:metrics-standard-projective-frames}
Let \(X=\mathbf P^n_K\) with homogeneous coordinates \(Z_0,\ldots,Z_n\) and
\(L=\mathcal O(1)\). The sections \(Z_i\) are global, and on the analytic
domain \(V_i=\{|Z_i|=\max_j|Z_j|\}\) the section \(Z_i\) generates \(L\).
Declare
\[
 \|Z_i(x)\|_{\mathrm{std}}
 =\frac{|Z_i(x)|}{\max_j|Z_j(x)|}
 \qquad\text{for all }i,
\]
the quotient being independent of the choice of homogeneous coordinates for
\(x\). On \(V_i\) this gives \(\|Z_i\|_{\mathrm{std}}=1\), so \(Z_i\) is a
frame of norm one there; the weight in that frame is
\(\phi_{Z_i}=0\). On the overlap \(V_i\cap V_j\) the transition function is
\(Z_j/Z_i\), whose absolute value is \(1\), and indeed
\(\phi_{Z_j}=\phi_{Z_i}-\log|Z_j/Z_i|=0\), as
\eqref{eq:metrics-weight-cocycle} requires. In the affine frame \(Z_0\) over
the whole chart \(\{Z_0\ne0\}\), with \(z_k=Z_k/Z_0\),
\[
 \phi_{Z_0}=-\log\|Z_0\|_{\mathrm{std}}
 =\log\max\{1,|z_1|,\ldots,|z_n|\},
\]
which for \(n=1\) is exactly \eqref{eq:p1-standard-weight}. The two
computations differ by \(-\log|Z_i/Z_0|\) on the overlap, as they must.
\end{example}

\subsection*{Normalization box: signs, weights, and where they are recorded}

Because this chapter is the place where valuations, divisors, metrics, and
measures are first put into one formula, we collect the binding conventions
here. Each is proved or adopted in the indicated place, and each is recorded
in the concordance of \cref{app:normalizations}.

\medskip
\noindent
\begin{tabular}{@{}p{.36\textwidth}p{.32\textwidth}p{.24\textwidth}@{}}
\toprule
convention & formula & recorded at\\
\midrule
valuations from absolute values
 & \(v_x(f)=-\log|f(x)|\)
 & \eqref{eq:appJ-valuation-sign}\\
uniformizer, discrete case
 & \(v_K(\pi)=1\), \(|\pi|=\rme^{-1}\)
 & \eqref{eq:appJ-uniformizer}\\
additive metric weight
 & \(\|s\|_\phi=\rme^{-\phi}\|s\|_0\)
 & \eqref{eq:appJ-na-weight}\\
change of reference
 & \(\phi'=\phi-\chi\)
 & \Cref{prop:metrics-torsor}\,(2)\\
Fubini--Study term of a section of \(mL\)
 & \(m^{-1}(\log|s|+\lambda)\)
 & \eqref{eq:appJ-fs-term}\\
model function of a vertical divisor
 & \(\varphi_D=-\log|f_D|=v_x(f_D)\)
 & \eqref{eq:appJ-model-divisor}\\
special-fiber test
 & \(\varphi_{\mathcal X_s}=1\)
 & \eqref{eq:appJ-special-fibre-test}\\
mixed model measure
 & \(\sum_Eb_E(\mathcal L_1\cdots\mathcal L_n|_E)\delta_{x_E}\)
 & \eqref{eq:appJ-cl-measure}\\
total mass
 & \((L_1\cdots L_n)\)
 & \eqref{eq:appJ-cl-total-mass}\\
Monge--Ampère of one metric
 & \(\MA(\phi)=c_1(L,\|\cdot\|)^n\)
 & \eqref{eq:appJ-model-ma}\\
probability normalization
 & \((L^n)^{-1}\MA(\phi)\), only if \((L^n)>0\)
 & \eqref{eq:appJ-prob-ma}\\
complex comparison
 & \(|s|_h^2=\rme^{-\varphi}|s|_{h_0}^2\), weight \(\varphi/2\)
 & \eqref{eq:appJ-complex-weight}\\
\bottomrule
\end{tabular}
\medskip

\noindent Three consequences are worth isolating, since they are the
fastest available sign tests.
\begin{enumerate}
\item \emph{Effective vertical divisors give nonnegative functions.} If
\(D\ge0\) is vertical then \(f_D\) is a regular function near every center,
so \(v_x(f_D)\ge0\) and \(\varphi_D\ge0\). Equivalently, twisting by an
effective vertical divisor \emph{shrinks} norms.
\item \emph{The special fiber is the constant \(1\).} Since
\(\mathcal X_s=\operatorname{div}(\pi)\) and \(v_x(\pi)=v_K(\pi)=1\) for
every \(x\in X^{\mathrm{an}}\), one has \(\varphi_{\mathcal X_s}\equiv1\);
adding \(c\,\mathcal X_s\) to a divisor adds the constant \(c\) to its model
function and multiplies the metric by \(\rme^{-c}\).
\item \emph{Masses are intersection numbers, not densities.} The total mass
of a mixed model measure is \((L_1\cdots L_n)\) and is never normalized to
one unless the word \emph{probability} is written; see
\cref{sec:appJ-model-ma}.
\end{enumerate}
Any formula in this chapter can be checked against these three tests on the
running conic of \eqref{eq:running-conic}, and we shall do so repeatedly.

\section{Algebraic model metrics}
\label{sec:metrics-model}

Assume first that \(K\) is discretely valued. Let \(\mathcal X\) be a proper
flat \(R\)-model of \(X\). A line bundle \(\mathcal L\) on \(\mathcal X\)
whose generic fiber is \(L\) defines a metric as follows. If \(e\) is a local
frame of \(\mathcal L\), declare
\[
                              \|e\|_{\mathcal L}=1
\]
on the corresponding analytic domain.

The phrase ``the corresponding analytic domain'' has to be taken seriously,
and it is the point at which the geometry of \cref{chap:berkovich} enters.
Let \(\mathcal U=\operatorname{Spec}B\subset\mathcal X\) be an affine open
subset on which \(\mathcal L\) is free, with frame \(e\). By properness of
\(\mathcal X\) over \(R\), every \(x\in X^{\mathrm{an}}\) has a unique center
\(c_{\mathcal X}(x)\in\mathcal X_s\) (\Cref{prop:centers-proper-separated}),
and by \Cref{prop:berk-reduction-anticontinuous} the set
\[
\begin{split}
 X^{\mathrm{an}}_{\mathcal U}
 &=\{x\in X^{\mathrm{an}}:c_{\mathcal X}(x)\in\mathcal U_s\}\\
 &=\{x\in X^{\mathrm{an}}:\rho_X(x)\in\mathcal U_K,\
   |f(x)|\le1\text{ for all }f\in B\}
\end{split}
\]
is a compact analytic domain of \(X^{\mathrm{an}}\), open in the special
fiber's Zariski topology in the anticontinuous sense. Since finitely many
affine opens cover \(\mathcal X\), finitely many such domains cover
\(X^{\mathrm{an}}\). It is on \(X^{\mathrm{an}}_{\mathcal U}\), and not on
the analytification of \(\mathcal U_K\), that \(e\) is declared to have norm
one.

More generally, suppose that \(m\geq1\) and
\(\mathcal L|_X\simeq L^{\otimes m}\). For a local section \(s\) of \(L\),
write \(s^{\otimes m}=fe\). Set
\[
                         \|s(x)\|_{\mathcal L}=|f(x)|^{1/m}.
\]
Changing \(e\) multiplies \(f\) by a model unit, whose norm is one. The
definition is therefore independent of all choices.

\begin{definition}
A metric obtained from a pair \((\mathcal X,m^{-1}\mathcal L)\) as above is
an \emph{algebraic model metric}, or simply a \emph{model metric}.
\end{definition}

The use of \(m^{-1}\mathcal L\) is shorthand for a rational model line
bundle; no root line bundle on \(\mathcal X\) is asserted to exist.

The two clauses of the construction---independence of the frame and
continuity---deserve complete proofs, since every later computation rests on
them and since the argument shows exactly which property of the center is
used.

\begin{proposition}
\label{prop:metrics-model-continuous}
Let \(K\) be discretely valued, let \(\mathcal X\) be a proper flat
\(R\)-model of the proper \(K\)-scheme \(X\), and let \(\mathcal L\) be a
line bundle on \(\mathcal X\) with \(\mathcal L|_X\simeq L^{\otimes m}\),
\(m\ge1\). Then:
\begin{enumerate}
\item the recipe above assigns to each \(x\in X^{\mathrm{an}}\) a norm on
\(L(x)\), independent of the chosen affine open, of the chosen frame of
\(\mathcal L\), and of the chosen presentation \(s^{\otimes m}=fe\);
\item the resulting metric \(\|\cdot\|_{\mathcal L}\) is continuous;
\item if \(\rho:\mathcal X'\to\mathcal X\) is a morphism of models
restricting to the identity on \(X\), then \(\rho^*\mathcal L\) induces the
same metric as \(\mathcal L\);
\item the assignment is multiplicative: \((\mathcal X,\mathcal L)\) and
\((\mathcal X,\mathcal L')\) induce metrics whose tensor product is the
metric induced by \(\mathcal L\otimes\mathcal L'\).
\end{enumerate}
\end{proposition}

\begin{proof}
(1) Let \(x\in X^{\mathrm{an}}\) with center \(\xi=c_{\mathcal X}(x)\), and
let \(e,e'\) be frames of \(\mathcal L\) on affine opens
\(\mathcal U,\mathcal U'\) both containing \(\xi\). On
\(\mathcal U\cap\mathcal U'\) we have \(e'=ue\) with
\(u\in\mathcal O_{\mathcal X}(\mathcal U\cap\mathcal U')^\times\); in
particular \(u\) and \(u^{-1}\) both lie in the local ring
\(\mathcal O_{\mathcal X,\xi}\). The center of \(x\) is by definition the
closed point of the image of \(\operatorname{Spec}R_x\to\mathcal X\), where
\(R_x=\{g:|g(x)|\le1\}\subset\mathcal H(x)\) is the valuation ring of \(x\);
hence \(\mathcal O_{\mathcal X,\xi}\) maps into \(R_x\), and therefore
\(|u(x)|\le1\) and \(|u^{-1}(x)|\le1\). Multiplicativity gives
\(|u(x)|=1\). Writing \(s^{\otimes m}=fe=f'e'\) we get \(f=uf'\) and thus
\(|f(x)|^{1/m}=|f'(x)|^{1/m}\). Finally, if \(s\ne0\) at \(x\) then \(f\) is
determined by \(s\) and \(e\), so no further choice is involved; and the
formula is homogeneous of degree one in \(s\), because replacing \(s\) by
\(gs\) replaces \(f\) by \(g^mf\) and \(|g^m|^{1/m}=|g|\).

(2) Continuity is local on the finite closed cover
\(\{X^{\mathrm{an}}_{\mathcal U}\}\) recalled above, \(\mathcal U\) running
over a finite affine cover of \(\mathcal X\) trivializing \(\mathcal L\). A
function on a topological space whose restrictions to the members of a
finite closed cover are continuous is continuous. On
\(X^{\mathrm{an}}_{\mathcal U}\) the frame \(e\) has constant norm one, and
for a local section \(s\) of \(L\) the function \(x\mapsto|f(x)|^{1/m}\) is
continuous because \(f\) is an analytic function there and \(t\mapsto
t^{1/m}\) is continuous. Hence \(\|\cdot\|_{\mathcal L}\) is continuous in
the frame \(e\), and therefore in every frame.

(3) Let \(x\in X^{\mathrm{an}}\) and let \(\xi'=c_{\mathcal X'}(x)\); then
\(\rho(\xi')=c_{\mathcal X}(x)=\xi\) because the center is the image of the
closed point of \(\operatorname{Spec}R_x\) and \(\rho\) is a morphism over
\(R\). If \(e\) is a frame of \(\mathcal L\) near \(\xi\), then
\(\rho^*e\) is a frame of \(\rho^*\mathcal L\) near \(\xi'\), and the
coefficient function transforms by \(f\mapsto f\circ\rho\), whose value at
\(x\) is \(f(x)\) by functoriality of evaluation
(\Cref{thm:berk-universal-property}). The two norms therefore agree.

(4) Immediate from the definitions: if \(e,e'\) are frames of
\(\mathcal L,\mathcal L'\) near the center, then \(e\otimes e'\) is a frame
of \(\mathcal L\otimes\mathcal L'\) there, and coefficients multiply.
\end{proof}

Assertion (3) is the reason models may always be replaced by dominating
ones, which by \Cref{prop:common-domination} exist for any two models. It
also shows that the metric depends on \(\mathcal L\) only through its class
in the direct limit over the system of models; the systematic formulation of
that limit is the Cartier b-divisor formalism of
\cref{app:intersection-bdivisors}.

\begin{proposition}\label{prop:metric-difference-model}
Let two model metrics on \(L\) be determined on one normal model
\(\mathcal X\). Choose a common tensor power \(M\) and let \(D_M\) be the
vertical Cartier divisor representing the quotient of the two resulting
extensions of \(L^{\otimes M}\).  Set \(D=M^{-1}D_M\).  Then
\[
                      \|\cdot\|_2=\|\cdot\|_1\rme^{-\varphi_D}.
\]
Here, if \(D_M=\operatorname{div}(f)\) locally,
\[
       \varphi_D=M^{-1}v_x(f)=-M^{-1}\log|f(x)|.
\]
\end{proposition}

\begin{proof}
After taking the common tensor power \(M\), let \(\mathcal L_1\) and
\(\mathcal L_2\) be extensions of \(L^{\otimes M}\). The
line bundle \(\mathcal L_2\otimes\mathcal L_1^{-1}\) is generically trivial.
A rational trivializing section has a vertical Cartier divisor \(D_M\), so
the line bundle is \(\mathcal O_{\mathcal X}(D_M)\).

If \(D_M=\operatorname{div}(f)\), then \(f^{-1}\) is a model frame of
\(\mathcal O_{\mathcal X}(D_M)\). Under the generic trivialization, the
original frame is \(f(f^{-1})\) and has norm \(|f|\). Hence the relative
additive weight on \(L^{\otimes M}\) is \(-\log|f|=v_x(f)\).  Dividing by
\(M\) when taking the induced \(M\)-th-root metrics on \(L\) gives
\(M^{-1}v_x(f)=\varphi_D\), as claimed.
\end{proof}

The proof deserves one comment, since the direction of the twist is the
commonest source of sign errors in the subject. At the integral level, the
line bundle \(\mathcal O_{\mathcal X}(D_M)\) is, by definition, the subsheaf
of the constant sheaf \(K(\mathcal X)\) generated locally by \(f^{-1}\) when
\(D_M=\operatorname{div}(f)\). Its distinguished frame is therefore
\(f^{-1}\), not \(f\); declaring \(\|f^{-1}\|=1\) forces
\(\|1\|=|f|\), that is, the weight \(-\log|f|=v_x(f)\) on the common tensor
power and the weight \(M^{-1}v_x(f)=\varphi_D\) on \(L\). Effective \(D\)
makes \(\varphi_D\ge0\) and the twisted metric smaller. This is the
convention of \eqref{eq:appJ-model-divisor}, and it is the one that makes
\(\varphi_{\mathcal X_s}=+1\).

The two sign tests from \cref{chap:skeleta} become metric tests:
\[
 \varphi_{\mathcal X_s}=v_x(\pi)=1,\qquad
 \|\cdot\|_{\mathcal L(\mathcal X_s)}
      =\rme^{-1}\|\cdot\|_{\mathcal L}.
\]
For a principal vertical divisor \(\operatorname{div}(h)\), the relative
weight is exactly \(v_x(h)\). These identities fix both the divisor sign and
the convention \(\|\cdot\|=\|\cdot\|_0\rme^{-\phi}\).

\subsection{Vertical divisors and model functions}
\label{sec:metrics-divisor-dictionary}

\Cref{prop:metric-difference-model} is the special case \(L=\mathcal O_X\)
of the whole theory, and in that case the metric is a function. We now
record the resulting dictionary in the form used throughout the book,
assembling the pieces proved in \cref{chap:models,chap:skeleta} and adding
the normalizations that were not needed there.

Write \(\mathcal D(X)\) for the set of continuous functions on
\(X^{\mathrm{an}}\) of the form \(\varphi_D\) for some vertical Cartier
divisor \(D\) on some proper flat model \(\mathcal X\) of \(X\), and
\(\mathcal D(X)_{\mathbf Q}\) for the corresponding set with \(D\) a
vertical \(\mathbf Q\)-Cartier divisor; these are the \emph{model
functions}, respectively \emph{rational model functions}, of \(X\).

\begin{proposition}
\label{prop:metrics-divisor-model-function}
Let \(K\) be discretely valued with \(v_K(\pi)=1\), and let \(X\) be a
proper, geometrically connected \(K\)-scheme.
\begin{enumerate}
\item \(\mathcal D(X)\) is a subgroup and \(\mathcal D(X)_{\mathbf Q}\) a
\(\mathbf Q\)-vector subspace of \(C^0(X^{\mathrm{an}})\), and
\(\mathcal D(X)_{\mathbf Q}=\mathbf Q\otimes_{\mathbf Z}\mathcal D(X)\).
\item For a fixed normal model \(\mathcal X\), the map \(D\mapsto\varphi_D\)
is injective on vertical \(\mathbf Q\)-Cartier divisors, and
\(\varphi_D(x_E)=a_E/b_E\) if \(D=\sum_Ea_EE\) and
\(\mathcal X_s=\sum_Eb_EE\).
\item \(\varphi_{\mathcal X_s}\equiv1\); consequently
\(\mathcal D(X)_{\mathbf Q}\) contains the constant functions with values in
\(\mathbf Q\), and \(\varphi_{D+c\mathcal X_s}=\varphi_D+c\).
\item \(D\ge0\) if and only if \(\varphi_D\ge0\) on \(X^{\mathrm{an}}\), and
this holds if and only if \(\varphi_D(x_E)\ge0\) for every irreducible
component \(E\) of \(\mathcal X_s\), where \(\mathcal X\) is a normal model
determining \(D\).
\item If \(h\in K^\times\) and
\(D=\operatorname{div}_{\mathcal X}(h)\), then \(D\) is vertical and
\(\varphi_D\equiv v_K(h)\).  For a general \(h\in K(X)^\times\), the full
divisor \(\operatorname{div}_{\mathcal X}(h)\) has a horizontal part and
does not define a vertical model function in this proposition.
\item Every \(\varphi\in\mathcal D(X)_{\mathbf Q}\) is bounded, and
\[
 \min_E\varphi(x_E)\le\varphi\le\max_E\varphi(x_E)
\]
on \(X^{\mathrm{an}}\), the extrema being taken over the components of the
special fiber of a normal model determining \(\varphi\).
\end{enumerate}
\end{proposition}

\begin{proof}
(1) If \(\varphi_{D_1}\) and \(\varphi_{D_2}\) are model functions determined
on models \(\mathcal X_1,\mathcal X_2\), choose a model \(\mathcal X\)
dominating both (\Cref{prop:common-domination}) and pull back the divisors;
by \Cref{lem:model-function-well-defined} this does not change the
associated functions, and \(\varphi_{D_1}+\varphi_{D_2}=\varphi_{D_1+D_2}\)
by multiplicativity of local equations. The same argument gives closure
under \(-1\) and under multiplication by rational scalars after clearing
denominators.

(2) This is \Cref{prop:divisor-model-function-injective-mod-fibre}; the
displayed value is the computation \(\varphi_D(x_E)=v_{x_E}(f_D)=
b_E^{-1}\operatorname{ord}_E(f_D)=a_E/b_E\), using
\Cref{prop:component-valuation}.

(3) The special fiber is cut out by \(\pi\), so \(\varphi_{\mathcal X_s}
=v_x(\pi)=v_K(\pi)=1\) for every \(x\), the middle equality because every
\(x\in X^{\mathrm{an}}\) restricts to \(v_K\) on \(K\). Additivity from (1)
gives the rest.

(4) If \(D\ge0\), a local equation \(f_D\) is regular at the center of any
\(x\), hence lies in \(R_x\) and has \(v_x(f_D)\ge0\). Conversely, if
\(\varphi_D(x_E)\ge0\) for every component \(E\), then \(a_E\ge0\) by (2),
so \(D\ge0\) as a divisor on the normal model; and then \(\varphi_D\ge0\)
everywhere by the first implication. This equivalence, valuative effectivity,
is used constantly later.

(5) A global local equation is \(h\), and every analytic point restricts to
\(v_K\) on \(K\), so
\(\varphi_D(x)=v_x(h)=v_K(h)\).  The warning about a general rational
function follows because its zeros and poles on the generic fibre close up
to horizontal components of its divisor.

(6) Let \(\mathcal X\) be a normal model determining \(D\) and set
\(c^+=\max_E\varphi_D(x_E)\). By (3) the divisor
\(c^+\mathcal X_s-D\) has model function \(c^+-\varphi_D\), whose values at
the \(x_E\) are \(c^+-a_E/b_E\ge0\); by (4) it is effective and its model
function is nonnegative, that is, \(\varphi_D\le c^+\). The lower bound is
symmetric. Boundedness also follows from continuity on the compact space
\(X^{\mathrm{an}}\), but the displayed bound is sharper and is the form used
in later compactness arguments.
\end{proof}

Assertion (6) is the reason model functions can be truncated and rescaled
freely: a model function is squeezed between the extreme values it takes at
the finitely many divisorial points of a determining model. Assertion (4)
is the valuative criterion for effectivity, and it will reappear as the
statement that a psh model function is detected on divisorial points.

\begin{example}[The standard model of the projective line]
\label{ex:metrics-p1-standard-model}
Let \(K\) be discretely valued, \(X=\mathbf P^1_K\), and
\(\mathcal X=\mathbf P^1_R\) with \(\mathcal L=\mathcal O_{\mathbf P^1_R}(1)\)
and \(L=\mathcal O(1)\). The special fiber is the reduced irreducible curve
\(\mathbf P^1_{\widetilde K}\), so \(b=1\) and there is exactly one
divisorial point of this model: the Gauss point \(\zeta_{0,1}\) of the unit
disc, which is the monomial valuation with \(v(T)=0\) in the coordinate
\(T=Z_1/Z_0\).

The induced metric is the standard one of
\cref{ex:metrics-standard-projective-frames}: on the analytic domain
\(\{|T|\le1\}\), which is exactly the set of points whose center lies in the
affine chart \(\operatorname{Spec}R[T]\) of the model, the section \(Z_0\) is
a frame of \(\mathcal L\), and its norm is one. Thus
\[
 -\log\|Z_0\|=\log\max\{1,|T|\},
\]
which is \eqref{eq:p1-standard-weight}; along the skeleton parametrized by
\(t=v(T)\) it is the function \(F(t)=\max\{0,-t\}\). The associated measure,
computed in \cref{sec:metrics-cl-measure}, is
\(\delta_{\zeta_{0,1}}\), of total mass \(1=\deg\mathcal O(1)\); this is
\eqref{eq:p1-standard-measure}.
\end{example}

\begin{example}[One blowup: the model function and its Laplacian]
\label{ex:metrics-p1-blowup}
\index{model function!blowup example}\index{Laplacian!and intersection numbers}
Continue with \(X=\mathbf P^1_K\) and let
\(\rho:\mathcal X'\to\mathbf P^1_R\) be the blowup of the closed point
\(p\in\mathbf P^1_{\widetilde K}\) given by \(T\equiv0\), that is, the
blowup of the ideal \(\mathfrak a=(\pi,T)\). Then \(\mathcal X'\) is a
regular projective model whose special fiber is
\[
 \mathcal X'_s=\widetilde C+E,
\]
where \(\widetilde C\) is the strict transform of
\(\mathbf P^1_{\widetilde K}\) and \(E\simeq\mathbf P^1_{\widetilde K}\) is
the exceptional divisor; both multiplicities are one, since
\(\operatorname{ord}_E(\pi)=1\). The two divisorial points are
\(x_{\widetilde C}=\zeta_{0,1}\) and \(x_E\), the latter being the monomial
valuation with
\[
 v_{x_E}(T)=1,\qquad v_{x_E}(\pi)=1,
\]
that is, the Gauss point \(\zeta_{0,|\pi|}\) of the disc \(|T|\le|\pi|\).

\emph{The model function.} Since \(\mathfrak a\cdot\mathcal O_{\mathcal X'}
=\mathcal O_{\mathcal X'}(-E)\), the function
\(\log|\mathfrak a|=-\varphi_E\) is computed by
\(\log|\mathfrak a|(x)=\max\{\log|\pi(x)|,\log|T(x)|\}\) at every point
whose center lies in the chart \(\operatorname{Spec}R[T]\), and
\(\log|\mathfrak a|=0\) elsewhere, because \(\mathfrak a\) is the unit ideal
away from \(p\). Hence
\begin{equation}
\label{eq:metrics-p1-blowup-hat}
 \varphi_E(x)=\min\{1,v_x(T)\}\ \text{ if }v_x(T)\ge0,
 \qquad
 \varphi_E(x)=0\ \text{ otherwise.}
\end{equation}
On the standard apartment, parametrized by \(t=v(T)\in[-\infty,+\infty]\),
this is the ``hat'' function \(\min\{1,t^+\}\): it vanishes for \(t\le0\),
rises with slope \(1\) on \([0,1]\), and is constantly \(1\) on the subtree
\(\{v_x(T)\ge1\}\), which is the closed disc of radius \(|\pi|\). The sign
tests hold: \(E\) is effective and \(\varphi_E\ge0\), with
\(\varphi_E(x_E)=1=a_E/b_E\) and \(\varphi_E(\zeta_{0,1})=0\), as
\Cref{prop:metrics-divisor-model-function}\,(2) requires.

\emph{Its Laplacian.} Apply the graph Laplacian
\eqref{eq:graph-laplacian}, which is the sum of outgoing slopes. At
\(\zeta_{0,1}\), the function \(\varphi_E\) has outgoing slope \(1\) in the
single direction pointing toward \(0\) and slope \(0\) in every other
direction, so the contribution is \(+1\). At \(x_E\), the outgoing slope in
the direction of \(\zeta_{0,1}\) is \(-1\), and \(\varphi_E\) is constant on
every other direction at \(x_E\), so the contribution is \(-1\). No other
point contributes. Hence
\begin{equation}
\label{eq:metrics-p1-blowup-laplacian}
 \Delta\varphi_E=\delta_{\zeta_{0,1}}-\delta_{x_E},
\end{equation}
of total mass \(0\), in accordance with \eqref{eq:lap-mass-zero}.

\emph{The twisted metric.} Consider the model line bundle
\(\mathcal L'=\rho^*\mathcal O(1)\otimes\mathcal O_{\mathcal X'}(-E)\), whose
generic fiber is again \(\mathcal O(1)\). By
\Cref{prop:metric-difference-model} its metric is the standard one twisted
by \(\varphi_{-E}=-\varphi_E\), so its weight in the frame \(Z_0\) is
\[
 \phi'=\log\max\{1,|T|\}-\min\{1,v(T)^+\}
 =\max\{-1,\log|T|\},
\]
the last equality being checked separately on \(v(T)\le0\), on
\(0\le v(T)\le1\), and on \(v(T)\ge1\). Thus the twisted metric is the
Fubini--Study metric attached to the two sections \(\pi Z_0\) and \(Z_1\)
of \(\mathcal O(1)\); see \cref{sec:metrics-fs}. Its curvature measure is,
by \eqref{eq:graph-curvature} and
\eqref{eq:metrics-p1-blowup-laplacian},
\[
 \delta_{\zeta_{0,1}}+\Delta(-\varphi_E)
 =\delta_{\zeta_{0,1}}-\delta_{\zeta_{0,1}}+\delta_{x_E}
 =\delta_{x_E},
\]
still of total mass \(1\). The same answer will be obtained in
\cref{ex:metrics-cl-p1-blowup} from intersection numbers on
\(\mathcal X'_s\), and the agreement of the two computations is the
one-dimensional case of the comparison between Laplacians and intersection
theory.
\end{example}

\section{Formal and piecewise linear metrics}
\label{sec:metrics-formal-pl}

Over a general complete nontrivially valued field, the valuation ring need
not be noetherian. Admissible formal models are then more flexible than
algebraic models. A line bundle \(\mathfrak L\) on an admissible formal
\(K^\circ\)-model \(\mathfrak X\) defines a \emph{formal metric} by declaring
formal frames to have norm one. Rational powers give
\(\mathbf Q\)-formal metrics.

The reader who has met only the algebraic construction may wonder what is
gained. Three things. First, admissibility---flatness together with
topological finite presentation---is the correct finiteness condition over a
valuation ring of arbitrary rank, and it is stable under the formal blowups
that generate the whole theory \cite{BoschLutkebohmert1993}. Second, the
generic fiber of a formal model is the compact analytic space itself, so no
separate properness argument is needed to know that centers exist. Third,
Raynaud's theorem makes admissible formal models and their blowups into a
cofiltered system with the same combinatorics as the algebraic one, so all
the domination arguments of \cref{chap:models} transpose. What is lost is
intersection theory: over a valuation ring of infinite rank there is no
Chow-theoretic degree in the form used in \cref{sec:metrics-cl-measure}, and
this is precisely why that section reverts to the discretely valued case.

There is also an intrinsic analytic definition.

\begin{definition}[Piecewise linear metric]
Let \(V\) be a paracompact strictly \(K\)-analytic space. A metric on a line
bundle \(L\) is \emph{piecewise linear} if \(V\) has a locally finite
analytic-domain covering \(V_i\) and frames \(s_i\) of \(L|_{V_i}\) such
that
\[
                              \|s_i\|\equiv1.
\]
It is \emph{piecewise \(\mathbf Q\)-linear} if locally a positive tensor
power is piecewise linear.
\end{definition}

The name refers to the tropical expression of the transition functions. On
a skeleton, a piecewise \(\mathbf Q\)-linear metric has rational piecewise
affine weights. Concretely: if \(s_i\) and \(s_j\) are two frames of norm
one on overlapping domains, then \(|s_i/s_j|\equiv1\) on the overlap, so
after passing to a skeleton the weight of the metric changes between charts
by an integral affine function---the valuation of a monomial. The picture to
keep in mind is \cref{ex:metrics-p1-blowup}, where the weight
\(\max\{-1,\log|T|\}\) is affine on each of the two pieces
\(\{v(T)\le1\}\) and \(\{v(T)\ge1\}\) of the apartment, with a corner at the
divisorial point of the exceptional divisor.

Three warnings about the terminology. It is not asserted that a piecewise
linear metric is piecewise linear in any Euclidean chart; the phrase refers
to the tropical structure just described. It is not asserted that the
covering can be taken finite, unless the space is compact. And piecewise
\(\mathbf Q\)-linear metrics are \emph{not} closed under uniform limits: as
we shall see, their uniform closure is the whole space of continuous
metrics, which is the content of part (3) of the next theorem.

\begin{theorem}
\label{thm:gm-metric-comparison}
Let \(K\) be an arbitrary complete nontrivially valued field.
\begin{enumerate}
\item On a paracompact strictly \(K\)-analytic space, formal metrics are
exactly piecewise linear metrics, and the same statement holds after
\(\mathbf Q\)-linearization.
\item If \(X\) is a proper \(K\)-scheme, algebraic model metrics, formal
metrics, and piecewise linear metrics on \(L^{\mathrm{an}}\) agree.
\item Every continuous metric on \(L^{\mathrm{an}}\) is a uniform limit of
piecewise \(\mathbf Q\)-linear metrics.
\end{enumerate}
\end{theorem}
\begin{proof}[Foundational metric-comparison input]
The three assertions are imported from Gubler--Martin with their published
scope.  The equivalence between formal and piecewise linear metrics,
including \(\mathbf Q\)-linearization, is
\cite[Theorems~1.1--1.2 and Proposition~2.10]{GM19}.  For a proper
algebraic scheme, the algebraization comparison in
\cite[Theorem~1.1]{GM19} identifies these with algebraic model metrics.
Finally, the uniform density statement is
\cite[Proposition~2.15 and Theorem~2.17]{GM19}.  Applying those results to
the analytic line bundle \(L^{\mathrm{an}}\) gives (1), (2), and (3),
respectively.  The density assertion does not claim that the approximants
are semipositive.
\end{proof}

\section{Fubini--Study metrics}
\label{sec:metrics-fs}

Fubini--Study metrics package a finite collection of sections into a metric.
Let \(m\geq1\), and let \(s_1,\ldots,s_N\in H^0(X,mL)\) have no common zero.
Choose real weights \(\lambda_1,\ldots,\lambda_N\). If
\(s_i=f_i\tau^{\otimes m}\) in a local frame \(\tau\) of \(L\), set
\[
 \phi_{\mathrm{FS},\tau}(x)
     =\frac1m\max_{1\leq i\leq N}
                    \{\log|f_i(x)|+\lambda_i\}.
\]
Under \(\tau'=g\tau\), every \(f_i\) changes to \(g^{-m}f_i\), and therefore
\[
                 \phi_{\mathrm{FS},\tau'}
                 =\phi_{\mathrm{FS},\tau}-\log|g|.
\]
The weights define a continuous metric, called a
\emph{Fubini--Study metric}.

Let us verify this last assertion and record the elementary estimates that
make the construction usable. Note that the normalization is
\eqref{eq:appJ-fs-term}: a section of \(mL\) enters with the factor
\(m^{-1}\), and its weight \(\lambda_i\) is added to \(\log|s_i|\), not
subtracted.

\begin{proposition}
\label{prop:metrics-fs-basic}
Let \(L\) be a line bundle on the proper \(K\)-scheme \(X\), let
\(m\ge1\), let \(s_1,\ldots,s_N\in H^0(X,mL)\) have no common zero, and let
\(\lambda\in\mathbf R^N\).
\begin{enumerate}
\item The local formulas \(\phi_{\mathrm{FS},\tau}\) satisfy the cocycle
rule \eqref{eq:metrics-weight-cocycle} and hence define a continuous metric
\(\|\cdot\|_{\mathrm{FS}(s,\lambda)}\) on \(L^{\mathrm{an}}\). It is
characterized among all continuous metrics on \(L\) by the two conditions
\begin{equation}
\label{eq:metrics-fs-characterization}
 \|s_i(x)\|\le\rme^{-\lambda_i}\quad(1\le i\le N),
 \qquad
 \max_{1\le i\le N}\rme^{\lambda_i}\|s_i(x)\|=1,
\end{equation}
where \(\|\cdot\|\) denotes the induced metric on \(mL\).
\item For two weight vectors \(\lambda,\mu\in\mathbf R^N\),
\[
 d_\infty\bigl(\|\cdot\|_{\mathrm{FS}(s,\lambda)},
 \|\cdot\|_{\mathrm{FS}(s,\mu)}\bigr)
 \le\frac1m\max_{1\le i\le N}|\lambda_i-\mu_i| .
\]
\item Adding a constant \(c\) to every \(\lambda_i\) adds \(c/m\) to the
weight; permuting the sections, or enlarging the family by sections
\(s_{N+1},\ldots\) with weights \(\lambda_j=-\infty\), does not change the
metric.
\item If \(t_1,\ldots,t_{N'}\in H^0(X,m'L)\) is a second base-point-free
family with weights \(\mu\), then the products \(s_it_j\) with weights
\(\lambda_i+\mu_j\) form a base-point-free family in \(H^0(X,(m+m')L)\)
whose Fubini--Study weight is
\[
 \frac{m}{m+m'}\phi_{\mathrm{FS}(s,\lambda)}
 +\frac{m'}{m+m'}\phi_{\mathrm{FS}(t,\mu)}.
\]
\end{enumerate}
\end{proposition}

\begin{proof}
(1) In the frame \(\tau'=g\tau\) one has
\(s_i=f_i\tau^{\otimes m}=(g^{-m}f_i)\tau'^{\otimes m}\), so
\(\phi_{\mathrm{FS},\tau'}=\frac1m\max_i\{\log|g^{-m}f_i|+\lambda_i\}
=\phi_{\mathrm{FS},\tau}-\log|g|\), which is
\eqref{eq:metrics-weight-cocycle}. Continuity of
\(\phi_{\mathrm{FS},\tau}\) holds because each \(\log|f_i|\) is continuous
with values in \([-\infty,\infty)\) and the maximum of finitely many such
functions is continuous provided it is finite; finiteness is exactly the
absence of a common zero, since \(f_i(x)=0\) for all \(i\) would force
\(s_i(\rho_X(x))=0\) for all \(i\). Now apply
\Cref{lem:metrics-weight-cocycle}. For
\eqref{eq:metrics-fs-characterization}, the induced metric on \(mL\) has
weight \(m\phi_{\mathrm{FS},\tau}\) in the frame \(\tau^{\otimes m}\), so
\[
 \rme^{\lambda_i}\|s_i(x)\|
 =\exp\Bigl(\lambda_i+\log|f_i(x)|
   -\max_j\{\log|f_j(x)|+\lambda_j\}\Bigr)\le1,
\]
with equality precisely at the indices realizing the maximum; since some
index always does, the maximum over \(i\) equals \(1\). Conversely, the two
conditions determine \(\|\cdot\|\) at every point, because the second one
prescribes the norm of at least one nonvanishing section.

(2) For real numbers \(a_i\) and weights \(\lambda,\mu\),
\[
 \max_i\{a_i+\lambda_i\}\le\max_i\{a_i+\mu_i\}
   +\max_i\{\lambda_i-\mu_i\},
\]
and symmetrically. Divide by \(m\) and take suprema over
\(X^{\mathrm{an}}\), using \eqref{eq:metrics-dinfty-sup}.

(3) The first two assertions are immediate from the formula. For the third,
a term with weight \(-\infty\) never contributes to the maximum; this
convention is used only to allow uniform notation and is never combined with
subtraction, in accordance with the extended-real rules of
\cref{chap:topological-tools}.

(4) The products have no common zero because a common zero of all
\(s_it_j\) would be a common zero of the \(s_i\) or of the \(t_j\). In a
frame,
\(\max_{i,j}\{\log|f_ig_j|+\lambda_i+\mu_j\}
=\max_i\{\log|f_i|+\lambda_i\}+\max_j\{\log|g_j|+\mu_j\}\),
since the maximum of a sum of independent families splits. Divide by
\(m+m'\), and use that the two maxima are respectively
\(m\phi_{\mathrm{FS}(s,\lambda)}\) and
\(m'\phi_{\mathrm{FS}(t,\mu)}\).
\end{proof}

\begin{example}[Projective space]
On \(\mathbf P^n_K\) with \(L=\mathcal O(1)\), the coordinate sections and
weights \(0\) give
\[
          \phi_{\mathrm{FS}}([Z_0:\cdots:Z_n])
             =\max_i\log|Z_i|
\]
in homogeneous notation. On the chart \(Z_0=1\), this becomes
\[
 \max\{0,\log|z_1|,\ldots,\log|z_n|\}.
\]
Its corner locus is the tropical
shadow of the standard projective model.
\end{example}

If the weights lie in the divisible value group of \(K\), clearing
denominators produces a lattice of sections and a relatively globally
generated model line bundle; the resulting Fubini--Study metric is a model
metric. Arbitrary real weights are uniform limits of such weights and give
continuous semipositive metrics when the valuation is nontrivial. We prove
both assertions, since they are the bridge between the section-theoretic and
the model-theoretic descriptions and since the second is the only place in
this chapter where the nontriviality of the valuation is used in an
essential way.

\begin{proposition}
\label{prop:metrics-fs-model}
Let \(K\) be discretely valued, let \(X\) be a projective \(K\)-scheme, let
\(s_1,\ldots,s_N\in H^0(X,mL)\) have no common zero, and suppose that every
\(\lambda_i\) lies in \(\Gamma_{\mathbf Q}:=\mathbf Q\otimes v(K^\times)\).
Then \(\|\cdot\|_{\mathrm{FS}(s,\lambda)}\) is an algebraic model metric,
and it is semipositive in the sense of
\cref{sec:metrics-semipositive}.
\end{proposition}

\begin{proof}
\emph{Step 1: reduction to integral weights in the value group.} Choose
\(q\ge1\) with \(q\lambda_i\in v(K^\times)\) for every \(i\). The family
\(s_i^{\otimes q}\in H^0(X,qmL)\) has no common zero, and in a frame
\[
 \frac1{qm}\max_i\bigl\{\log|f_i^q|+q\lambda_i\bigr\}
 =\frac1m\max_i\bigl\{\log|f_i|+\lambda_i\bigr\},
\]
so the two data define the same metric. We may therefore assume
\(\lambda_i\in v(K^\times)\) for all \(i\).

\emph{Step 2: absorbing the weights into the sections.} Since
\(v(K^\times)\) is a group, for each \(i\) there is \(c_i\in K^\times\) with
\(\log|c_i|=\lambda_i\). Replacing \(s_i\) by \(t_i:=c_is_i\) replaces
\(\log|f_i|\) by \(\log|f_i|+\lambda_i\), so
\(\phi_{\mathrm{FS}(s,\lambda)}=\phi_{\mathrm{FS}(t,0)}\). The family
\((t_i)\) still has no common zero. We may therefore assume
\(\lambda=0\).

\emph{Step 3: the associated morphism.} Base-point-freeness of
\((t_i)_{1\le i\le N}\) gives a morphism
\[
 \psi:X\longrightarrow\mathbf P^{N-1}_K,
 \qquad \psi^*\mathcal O(1)\simeq mL,\qquad \psi^*Z_i=t_i .
\]
By \Cref{ex:metrics-standard-projective-frames}, the standard metric on
\(\mathcal O(1)\) is the model metric of
\(\mathcal O_{\mathbf P^{N-1}_R}(1)\) on \(\mathbf P^{N-1}_R\). Fix a point
of \(X^{\mathrm{an}}\) and choose an index \(j\) with \(t_j\ne0\) at that
point, which is possible by base-point-freeness; on the neighborhood where
\(t_j\ne0\), the section \(t_j=f_j\tau^{\otimes m}\) is a frame of \(mL\),
and the weight of the pullback metric in the frame \(\tau^{\otimes m}\) is
\[
 \psi^*\Bigl(\max_i\log|Z_i/Z_j|\Bigr)+\log|f_j|
 =\max_i\log|f_i|=m\,\phi_{\mathrm{FS}(t,0),\tau},
\]
where we used \(f_i/f_j=t_i/t_j=\psi^*(Z_i/Z_j)\) and the transformation
rule \eqref{eq:metrics-weight-cocycle}. Since the point was arbitrary and
both sides are continuous, the identity holds on all of
\(X^{\mathrm{an}}\). Thus
\(\|\cdot\|_{\mathrm{FS}(t,0)}\) is the \(m\)-th root of the pullback under
\(\psi\) of the standard metric.

\emph{Step 4: a model computing the pullback.} Choose any projective
\(R\)-model \(\mathcal X_0\) of \(X\); one exists by taking the flat closure
of \(X\) in a projective space over \(R\)
(\Cref{lem:flat-closure-generic-fibre}). Let \(\mathcal X\) be the flat
closure of the graph of \(\psi\) inside
\(\mathcal X_0\times_R\mathbf P^{N-1}_R\); it is a projective \(R\)-model of
\(X\), and the second projection gives a morphism
\(\Psi:\mathcal X\to\mathbf P^{N-1}_R\) extending \(\psi\). Put
\(\mathcal L=\Psi^*\mathcal O_{\mathbf P^{N-1}_R}(1)\), a line bundle on
\(\mathcal X\) with \(\mathcal L|_X\simeq mL\). By
\Cref{prop:metrics-model-continuous}\,(3) and functoriality of centers, the
model metric of \((\mathcal X,m^{-1}\mathcal L)\) is the \(m\)-th root of the
pullback of the standard model metric, hence equals
\(\|\cdot\|_{\mathrm{FS}(t,0)}\) by Step 3.

\emph{Step 5: semipositivity.} Let \(C\subset\mathcal X_s\) be an integral
proper curve. Then
\[
 \deg_{\mathcal L}(C)=\deg\bigl(\Psi|_C^*\mathcal O(1)\bigr).
\]
If \(\Psi(C)\) is a point, the restriction is trivial and the degree is
\(0\). Otherwise \(\Psi(C)\) is an integral curve in
\(\mathbf P^{N-1}_{\widetilde K}\), the morphism \(C\to\Psi(C)\) is
finite, and the projection formula gives
\[
 \deg_{\mathcal L}(C)
 =\deg(C/\Psi(C))\cdot\deg\bigl(\mathcal O(1)|_{\Psi(C)}\bigr)\ge0,
\]
because \(\mathcal O(1)\) is ample on
\(\mathbf P^{N-1}_{\widetilde K}\). Hence
\(\mathcal L\) is relatively nef and the metric is semipositive.
\end{proof}

\begin{proposition}
\label{prop:metrics-fs-real-weights}
Let \(K\) be discretely valued and let \(X\), \(L\), \(m\), and
\((s_i)\) be as in \Cref{prop:metrics-fs-model}, but with arbitrary
\(\lambda\in\mathbf R^N\). Then
\(\|\cdot\|_{\mathrm{FS}(s,\lambda)}\) is a uniform limit of Fubini--Study
metrics with weights in \(\Gamma_{\mathbf Q}\), hence a continuous
semipositive metric.
\end{proposition}

\begin{proof}
Since \(K\) is nontrivially valued, \(v(K^\times)\) is a nonzero subgroup of
\(\mathbf R\), so \(\Gamma_{\mathbf Q}=\mathbf Q\otimes v(K^\times)\) is a
nonzero \(\mathbf Q\)-vector subspace of \(\mathbf R\) and is therefore
dense. Choose \(\lambda^{(j)}\in\Gamma_{\mathbf Q}^N\) with
\(\max_i|\lambda_i^{(j)}-\lambda_i|\le1/j\). By
\Cref{prop:metrics-fs-basic}\,(2),
\[
 d_\infty\bigl(\|\cdot\|_{\mathrm{FS}(s,\lambda^{(j)})},
 \|\cdot\|_{\mathrm{FS}(s,\lambda)}\bigr)\le\frac1{mj}
 \xrightarrow[j\to\infty]{}0 .
\]
Each approximant is a semipositive model metric by
\Cref{prop:metrics-fs-model}, so the limit is a continuous semipositive
metric by the definition in \cref{sec:metrics-semipositive}.
\end{proof}

Two remarks on the hypotheses. Discreteness of the valuation was used only
through the availability of algebraic models with the intersection-theoretic
vocabulary of \cref{chap:models}; over an arbitrary complete nontrivially
valued field the same statements hold with formal models, by
\Cref{thm:gm-metric-comparison}. Nontriviality was used, and cannot be
removed: over a trivially valued field \(v(K^\times)=0\), the argument of
\Cref{prop:metrics-fs-real-weights} breaks down completely, and
Fubini--Study functions with arbitrary real weights become primitive objects
rather than limits of model data. That is the starting point of
\cref{part:bj}.

\begin{example}[A Fubini--Study metric on \(\mathcal O(1)\), written out]
\label{ex:metrics-fs-p1-written}
\index{Fubini--Study metric!computed on the projective line}
Let \(K\) be discretely valued with \(v(\pi)=1\), let
\(X=\mathbf P^1_K\) with coordinates \([Z_0:Z_1]\), \(T=Z_1/Z_0\), and let
\(L=\mathcal O(1)\), \(m=1\). Take the two sections
\[
 s_1=\pi Z_0,\qquad s_2=Z_1,
\]
with weights \(\lambda=0\). They have no common zero. In the frame
\(\tau=Z_0\) one has \(f_1=\pi\) and \(f_2=T\), so
\[
 \phi_{\mathrm{FS},Z_0}
 =\max\{\log|\pi|,\log|T|\}
 =\max\{-1,\log|T|\} .
\]
Along the apartment parametrized by \(t=v(T)\), this is
\(F(t)=\max\{-1,-t\}\): affine of slope \(-1\) for \(t\le1\) and constant
\(-1\) for \(t\ge1\), with a single corner at \(t=1\). This is exactly the
weight computed in \cref{ex:metrics-p1-blowup} from the blowup of
\((\pi,T)\), and the two descriptions agree as they must: clearing
denominators in the Fubini--Study data produces the ideal \((\pi,T)\), whose
normalized blowup is the model \(\mathcal X'\) used there.

Its curvature measure is a single Dirac mass at the divisorial point of the
exceptional divisor,
\[
 c_1\bigl(\mathcal O(1),\|\cdot\|_{\mathrm{FS}}\bigr)=\delta_{x_E},
 \qquad x_E=\zeta_{0,|\pi|},
\]
of total mass \(1=\deg\mathcal O(1)\). Replacing \(s_1\) by \(\pi^aZ_0\)
with \(a\in\mathbf Z_{>0}\) moves the corner to \(t=a\) and moves the atom to
\(\zeta_{0,|\pi|^a}\).  More generally, assign the real weight
\(-\lambda\) to \(Z_0\) and weight zero to \(Z_1\).  The corner then lies at
\(t=\lambda\).  This is an integral model metric for
\(\lambda\in v(K^\times)\), a \(\mathbf Q\)-model metric for
\(\lambda\in\mathbf Qv(K^\times)\), and otherwise only a real weighted
Fubini--Study metric.  In the last case the corner is quasimonomial and
non-divisorial.
\end{example}

\begin{example}[Two corners: a Fubini--Study metric from three sections]
\label{ex:metrics-fs-two-corners}
Keep \(X=\mathbf P^1_K\) and \(L=\mathcal O(1)\), and take \(m=2\) with the
three sections of \(\mathcal O(2)\)
\[
 s_1=\pi^4Z_0^2,\qquad s_2=\pi Z_0Z_1,\qquad s_3=Z_1^2,
\]
all weights zero. They have no common zero, since \(s_1\) and \(s_3\) do
not vanish simultaneously. In the frame \(Z_0^{\otimes2}\) one has
\(f_1=\pi^4\), \(f_2=\pi T\), \(f_3=T^2\), so with \(t=v(T)\) and
\(\log|T|=-t\),
\[
 F(t)=\tfrac12\max\{-4,\,-1-t,\,-2t\}
 =\begin{cases}
   -t,&t\le1,\\[1mm]
   \dfrac{-1-t}2,&1\le t\le3,\\[2mm]
   -2,&t\ge3 .
  \end{cases}
\]
Indeed \(-2t\ge-1-t\) exactly for \(t\le1\), and \(-1-t\ge-4\) exactly for
\(t\le3\); the middle section is the maximal one precisely on \([1,3]\),
which is nonempty because the three exponents \(-4,-1,0\) form a strictly
concave sequence. The function is convex and piecewise affine with slopes
\(-1,-\tfrac12,0\) and corners at \(t=1\) and \(t=3\).

To compute the curvature measure, subtract the standard weight
\(F_{\mathrm{std}}(t)=\max\{0,-t\}\) of
\cref{ex:metrics-p1-standard-model}. The relative weight is
\[
 \varphi=F-F_{\mathrm{std}}
 =\begin{cases}
   0,&t\le0,\\
   -t,&0\le t\le1,\\
   \tfrac{-1-t}2,&1\le t\le3,\\
   -2,&t\ge3,
  \end{cases}
\]
a model function, constant on every branch of the tree leaving the
apartment. Summing outgoing slopes as in \eqref{eq:graph-laplacian} gives
contributions \(-1\) at \(t=0\), \(1-\tfrac12=\tfrac12\) at \(t=1\), and
\(\tfrac12-0=\tfrac12\) at \(t=3\), so
\(\Delta\varphi=-\delta_{\zeta_{0,1}}+\tfrac12\delta_{t=1}
+\tfrac12\delta_{t=3}\), of total mass \(0\) as
\eqref{eq:lap-mass-zero} requires. Adding the curvature
\(\delta_{\zeta_{0,1}}\) of the standard metric, as in
\eqref{eq:graph-curvature}, yields
\[
 c_1\bigl(\mathcal O(1),\|\cdot\|_{\mathrm{FS}}\bigr)
 =\tfrac12\,\delta_{\zeta_{0,|\pi|}}+\tfrac12\,\delta_{\zeta_{0,|\pi|^3}},
\]
of total mass \(1=\deg\mathcal O(1)\). The two masses are the two slope
jumps. This metric is a model metric for the model obtained by blowing up
the ideal \((\pi^4,\pi T,T^2)\) and normalizing, and the half-integral
masses are the reason the definition of a model metric allows the rational
twist \(m^{-1}\mathcal L\): the underlying model line bundle extends
\(\mathcal O(2)\), not \(\mathcal O(1)\).
\end{example}

\begin{remark}
A model line bundle that is relatively globally generated gives a
Fubini--Study metric. A relatively nef line bundle need not be semiample, so
one must not identify all semipositive model metrics with individual
Fubini--Study metrics. Approximation results later recover the useful
relationship. Over a trivially valued field, real weighted Fubini--Study
functions are basic objects in their own right and are not defined by
silently applying the preceding value-group argument.
\end{remark}

\section{Model functions and the density theorem}
\label{sec:metrics-model-functions}

The metrics constructed so far form a small set inside all continuous
metrics: they are indexed by algebraic model data, whereas continuous metrics form an
infinite-dimensional Banach torsor by \Cref{lem:metrics-dinfty-complete}.
The whole theory nevertheless works, because the small set is uniformly
dense. This section proves that, in the following function-theoretic form:
the space \(\mathcal D(X)_{\mathbf Q}\) of rational model functions of
\cref{sec:metrics-divisor-dictionary} is dense in \(C^0(X^{\mathrm{an}})\).

The result is due to Gubler, who proved it in his work on local heights of
subvarieties over non-Archimedean fields; the version for piecewise linear
metrics over an arbitrary complete nontrivially valued field is
\cite[Theorem~2.17]{GM19}, already quoted as
\Cref{thm:gm-metric-comparison}\,(3), and the model-function version in the
discretely valued setting used by Boucksom--Favre--Jonsson is
\cite[\S2]{BFJ16}. Because every later approximation argument---the
extension of Chambert--Loir measures in \cref{sec:metrics-continuous-cl},
the Radon-measure duality of \cref{part:bfj}, the regularization theorems of
\cref{part:bj}---rests on it, we give a complete proof in the regime of this
chapter. The proof has three ingredients: model functions form a lattice,
they contain a dense set of constants, and they separate points. The first
and third are geometric and are proved below; the passage from these to
density is the lattice form of the Stone--Weierstrass theorem, which we also
prove, since its hypotheses have to be matched exactly.

\begin{lemma}
\label{lem:metrics-model-lattice}
Let \(K\) be discretely valued and let \(X\) be an integral proper
\(K\)-scheme. If \(\varphi_1,\varphi_2\in\mathcal D(X)_{\mathbf Q}\), then
\[
 \min\{\varphi_1,\varphi_2\}\in\mathcal D(X)_{\mathbf Q}
 \qquad\text{and}\qquad
 \max\{\varphi_1,\varphi_2\}\in\mathcal D(X)_{\mathbf Q}.
\]
\end{lemma}

\begin{proof}
Since \(\min\{q\varphi_1,q\varphi_2\}=q\min\{\varphi_1,\varphi_2\}\) for
\(q>0\), and since \(\mathcal D(X)_{\mathbf Q}\) is a \(\mathbf Q\)-vector
space by \Cref{prop:metrics-divisor-model-function}\,(1), we may clear
denominators and assume that \(\varphi_i=\varphi_{D_i}\) with \(D_i\)
vertical Cartier divisors determined on one integral proper model
\(\mathcal X\); a common model exists by \Cref{prop:common-domination}, and
pulling the divisors back does not change the functions by
\Cref{lem:model-function-well-defined}. Adding a large multiple of
\(\mathcal X_s\) to both, which by
\Cref{prop:metrics-divisor-model-function}\,(3) adds the same constant to
both functions and hence commutes with \(\min\), we may also assume that
\(D_1\) and \(D_2\) are effective.

Let \(\mathfrak a=\mathcal O_{\mathcal X}(-D_1)+\mathcal O_{\mathcal X}(-D_2)
\subset\mathcal O_{\mathcal X}\); locally, if \(D_i=\operatorname{div}(f_i)\)
with \(f_i\) regular, then \(\mathfrak a=(f_1,f_2)\). Two local
presentations differ by units of the local ring at the center of any given
analytic point, so, exactly as in the proof of
\Cref{prop:metrics-model-continuous}\,(1), the function
\[
 \log|\mathfrak a|(x):=\max\bigl\{\log|f_1(x)|,\log|f_2(x)|\bigr\}
 =-\min\{\varphi_{D_1}(x),\varphi_{D_2}(x)\}
\]
is well defined on \(X^{\mathrm{an}}\). Since \(D_1,D_2\) are vertical,
\(\mathfrak a\) restricts to the unit ideal on the generic fiber.

Let \(\rho:\mathcal X'\to\mathcal X\) be the blowup of \(\mathfrak a\).
It is projective, and it is an isomorphism over the generic fiber because
\(\mathfrak a\) is invertible there; hence \(\mathcal X'\) is integral,
dominates \(\operatorname{Spec}R\), and is therefore a proper flat model of
\(X\) by \Cref{lem:torsion-free-flat-valuation}. By construction
\(\mathfrak a\cdot\mathcal O_{\mathcal X'}=\mathcal O_{\mathcal X'}(-E)\) for
an effective Cartier divisor \(E\), which is vertical because
\(\mathfrak a\) is the unit ideal on the generic fiber.

We claim \(\log|\mathfrak a|=-\varphi_E\). Let \(x\in X^{\mathrm{an}}\) and
let \(\xi'=c_{\mathcal X'}(x)\), which lies over \(\xi=c_{\mathcal X}(x)\).
Near \(\xi'\), one of the pulled-back generators, say \(f_1\), generates
\(\mathfrak a\cdot\mathcal O_{\mathcal X'}\), so \(f_2/f_1\) is regular at
\(\xi'\) and therefore \(|f_2(x)|\le|f_1(x)|\) by the argument of
\Cref{prop:metrics-model-continuous}\,(1). Hence
\(\max\{\log|f_1(x)|,\log|f_2(x)|\}=\log|f_1(x)|=-\varphi_E(x)\), since
\(f_1\) is a local equation of \(E\) at \(\xi'\). This proves the claim, and
therefore
\[
 \min\{\varphi_{D_1},\varphi_{D_2}\}=\varphi_E
 \in\mathcal D(X).
\]
Finally \(\max\{\varphi_1,\varphi_2\}
=-\min\{-\varphi_1,-\varphi_2\}\), and
\(\mathcal D(X)_{\mathbf Q}\) is stable under negation.
\end{proof}

In multiplicative language, \Cref{lem:metrics-model-lattice} says that the
pointwise maximum of two model \emph{metrics}, that is, the pointwise
minimum of the corresponding norms, is again a model metric; the
corresponding statement about semipositivity is genuinely harder and is
treated in \cref{sec:metrics-semipositive}. Notice which ideal was blown up:
the sum of the two ideal sheaves, not their product or intersection. The
sum computes the maximum of the two functions \(\log|f_i|\), hence the
minimum of the two weights.

\begin{lemma}
\label{lem:metrics-model-separation}
Let \(K\) be complete, discretely valued, and nontrivially valued, and let
\(X\) be a projective \(K\)-scheme. For any two distinct points
\(x,y\in X^{\mathrm{an}}\), there is \(\varphi\in\mathcal D(X)_{\mathbf Q}\)
with \(\varphi(x)\ne\varphi(y)\).
\end{lemma}

\begin{proof}
\emph{Step 1: reduction to projective space.} Fix a closed immersion
\(\iota:X\hookrightarrow\mathbf P^N_K\). By
\cref{sec:berk-functoriality}, \(\iota^{\mathrm{an}}\) is a closed immersion
of analytic spaces, in particular injective, so
\(\iota^{\mathrm{an}}(x)\ne\iota^{\mathrm{an}}(y)\). Moreover the
restriction to \(X\) of a model function on \(\mathbf P^N_K\) is a model
function on \(X\): if \(\varphi=\varphi_D\) for a vertical divisor \(D\) on
a model \(\mathcal Y\) of \(\mathbf P^N_K\), let \(\mathcal X\) be the flat
closure of \(X\) in \(\mathcal Y\) (\Cref{lem:flat-closure-generic-fibre});
it is a model of \(X\), the pullback of \(D\) to \(\mathcal X\) is a
vertical Cartier divisor provided \(\mathcal X\not\subset\operatorname{Supp}
D\), which holds because \(\mathcal X\) is flat and \(D\) is vertical, and
its model function is \(\varphi\circ\iota^{\mathrm{an}}\) by functoriality of
evaluation. It therefore suffices to treat \(X=\mathbf P^N_K\).

\emph{Step 2: the separating family.} For a homogeneous polynomial \(F\) of
degree \(d\ge1\) in \(Z_0,\ldots,Z_N\) and \(c\in K^\times\), set
\[
 \psi_{F}(z)=\log|F(z)|-d\log\max_k|Z_k(z)|\in[-\infty,+\infty),
\]
\[
 \varphi_{F,c}=\max\bigl\{\psi_F+\log|c|,\,0\bigr\},
\]
where the quotient defining \(\psi_F\) is computed in any system of
homogeneous coordinates and is therefore well defined; it is continuous with
values in \([-\infty,+\infty)\), and \(\varphi_{F,c}\) is continuous and
real valued. We claim that
\(\varphi_{F,c}\in\mathcal D(\mathbf P^N_K)_{\mathbf Q}\). Indeed,
\(d^{-1}\varphi_{F,c}\) is the difference of the weights of two
Fubini--Study metrics on \(\mathcal O(1)\), namely the metric attached to
the base-point-free family \(\{cF,Z_0^d,\ldots,Z_N^d\}\subset
H^0(\mathbf P^N,\mathcal O(d))\) with zero weights, and the metric attached
to \(\{Z_0^d,\ldots,Z_N^d\}\); both families have no common zero because
the \(Z_k^d\) have none. Both metrics are model metrics by
\Cref{prop:metrics-fs-model}, so their difference is a model function by
\Cref{prop:metric-difference-model}, applied after passing to a common
model; multiplying by the integer \(d\) keeps us inside the
\(\mathbf Q\)-vector space \(\mathcal D(\mathbf P^N_K)_{\mathbf Q}\).

\emph{Step 3: separation in a common chart.} For \(0\le i\le N\) put
\(U_i=\{z:|Z_i(z)|=\max_k|Z_k(z)|\}\). These are compact analytic domains
covering \((\mathbf P^N_K)^{\mathrm{an}}\), and in the affine coordinates
\(w_k=Z_k/Z_i\) the domain \(U_i\) is the closed unit polydisc
\(\mathcal M(K\langle w\rangle)\) by
\Cref{lem:berk-polydisc-exhaustion}. Suppose first that \(x,y\in U_i\) for
some \(i\). Being distinct points of an affinoid spectrum, they differ on
some element \(f\in K\langle w\rangle\), say
\(\delta:=\bigl||f|_x-|f|_y\bigr|>0\). Polynomials are dense in
\(K\langle w\rangle\) for the Gauss norm, and every point \(z\) of the
polydisc satisfies \(\bigl||f|_z-|P|_z\bigr|\le|f-P|_z\le\|f-P\|\), by the
ultrametric inequality and the norm bound of
\Cref{lem:berk-spectral-bound}. Choosing a polynomial \(P\) with
\(\|f-P\|<\delta/2\) therefore gives \(|P|_x\ne|P|_y\). Write
\(P=F/Z_i^d\) with \(F\) homogeneous of degree \(d\ge1\); enlarging \(d\)
if necessary by multiplying \(F\) by a power of \(Z_i\), which changes
neither \(P\) nor the conclusion, we may assume \(d\ge1\). On \(U_i\) we
have \(\max_k|Z_k|=|Z_i|\), so \(\psi_F=\log|P|\) there, and hence
\(\psi_F(x)\ne\psi_F(y)\) in \([-\infty,\infty)\).

Choose \(c\in K^\times\) with \(\log|c|>0\) large; this is possible because
\(K\) is nontrivially valued. If both \(\psi_F(x)\) and \(\psi_F(y)\) are
finite, choose \(\log|c|>-\min\{\psi_F(x),\psi_F(y)\}\); then both
\(\psi_F+\log|c|\) are positive at \(x\) and \(y\), so
\(\varphi_{F,c}(x)-\varphi_{F,c}(y)=\psi_F(x)-\psi_F(y)\ne0\). If, say,
\(\psi_F(x)=-\infty\ne\psi_F(y)\), choose \(\log|c|>-\psi_F(y)\); then
\(\varphi_{F,c}(x)=0<\psi_F(y)+\log|c|=\varphi_{F,c}(y)\).

\emph{Step 4: separation in different charts.} Suppose no \(U_i\) contains
both points. Choose \(j\) with \(y\in U_j\); then \(x\notin U_j\), that is,
\(\psi_{Z_j}(x)<0=\psi_{Z_j}(y)\). Apply Step 3's construction with
\(F=Z_j\), \(d=1\), and any \(c\) with \(\log|c|>0\): then
\(\varphi_{Z_j,c}(y)=\log|c|\), while
\(\varphi_{Z_j,c}(x)=\max\{\psi_{Z_j}(x)+\log|c|,0\}<\log|c|\), because
\(\psi_{Z_j}(x)<0\) and \(\log|c|>0\). The two values differ.
\end{proof}

\begin{proposition}
\label{prop:metrics-model-function-density}
Let \(K\) be a complete, discretely valued, nontrivially valued field and let
\(X\) be an integral projective \(K\)-scheme. Then the space
\(\mathcal D(X)_{\mathbf Q}\) of rational model functions is uniformly dense
in \(C^0(X^{\mathrm{an}})\).
\end{proposition}

\begin{proof}
Write \(Z=X^{\mathrm{an}}\), a compact Hausdorff space by
\Cref{thm:berk-proper-compact}, and let
\(A=\mathcal D(X)_{\mathbf Q}\subset C^0(Z)\) and
\(\overline A\) its closure for the sup norm.

\emph{Step 1: \(\overline A\) is a real vector lattice containing the
constants.} By \Cref{prop:metrics-divisor-model-function}\,(1), \(A\) is a
\(\mathbf Q\)-vector space; its closure is stable under addition, and stable
under multiplication by arbitrary reals because
\(\|q_nf-rf\|_\infty=|q_n-r|\,\|f\|_\infty\to0\) for rationals
\(q_n\to r\). Hence \(\overline A\) is an \(\mathbf R\)-vector subspace of
\(C^0(Z)\). By \Cref{prop:metrics-divisor-model-function}\,(3) it contains
all rational constants, hence all real constants. It is stable under
\(\max\) and \(\min\): these operations preserve \(A\) by
\Cref{lem:metrics-model-lattice}, and they are nonexpanding for the sup
norm, since
\(\|\max\{f,g\}-\max\{f',g'\}\|_\infty\le
\max\{\|f-f'\|_\infty,\|g-g'\|_\infty\}\).

\emph{Step 2: two-point interpolation.} Let \(z\ne z'\) in \(Z\) and
\(\alpha,\beta\in\mathbf R\). By \Cref{lem:metrics-model-separation} there
is \(h\in A\) with \(h(z)\ne h(z')\). Then
\[
 f=\alpha+(\beta-\alpha)\frac{h-h(z)}{h(z')-h(z)}\in\overline A
\]
satisfies \(f(z)=\alpha\) and \(f(z')=\beta\), using Step 1.

\emph{Step 3: the lattice Stone--Weierstrass argument.} Let
\(g\in C^0(Z)\) and \(\varepsilon>0\). For each ordered pair
\((z,z')\) with \(z\ne z'\), choose by Step 2 a function
\(f_{z,z'}\in\overline A\) with \(f_{z,z'}(z)=g(z)\) and
\(f_{z,z'}(z')=g(z')\); for \(z=z'\) take the constant \(g(z)\). Fix
\(z\in Z\). For each \(z'\in Z\) the set
\[
 U_{z'}=\{w\in Z: f_{z,z'}(w)<g(w)+\varepsilon\}
\]
is open and contains \(z'\). By compactness there are
\(z'_1,\ldots,z'_k\) with \(Z=\bigcup_jU_{z'_j}\); put
\(f_z=\min_jf_{z,z'_j}\in\overline A\). Then \(f_z<g+\varepsilon\) on
\(Z\), and \(f_z(z)=g(z)\). Next, the sets
\(V_z=\{w:f_z(w)>g(w)-\varepsilon\}\) are open and contain \(z\); by
compactness there are \(z_1,\ldots,z_l\) with \(Z=\bigcup_iV_{z_i}\). Put
\(f=\max_if_{z_i}\in\overline A\). Then \(f>g-\varepsilon\) on \(Z\),
because every \(w\) lies in some \(V_{z_i}\), and \(f<g+\varepsilon\) on
\(Z\), because each \(f_{z_i}\) satisfies this bound. Hence
\(\|f-g\|_\infty\le\varepsilon\), so \(g\in\overline{\overline A}=\overline
A\).

Since \(\varepsilon\) and \(g\) were arbitrary,
\(\overline A=C^0(Z)\).
\end{proof}

Three comments on the scope of \Cref{prop:metrics-model-function-density}.

First, the projectivity hypothesis was used only in
\Cref{lem:metrics-model-separation}, through a projective embedding. For a
proper \(X\), and indeed over an arbitrary complete nontrivially valued
field, the corresponding statement for piecewise linear metrics is
\Cref{thm:gm-metric-comparison}\,(3); applying it to the trivial line bundle
and using part (2) of the same theorem recovers density of model functions
in \(C^0(X^{\mathrm{an}})\) in that generality. The point of the proof just
given is that in the discretely valued projective case it is elementary and
self-contained, and that it exhibits explicitly the functions doing the
approximating: maxima of finitely many Fubini--Study differences.

Second, the statement is about \emph{functions}. It becomes a statement
about metrics through \Cref{prop:metrics-torsor}: fixing one model metric on
\(L\), which exists by \Cref{prop:line-bundle-extension} after passing to a
positive tensor power, one obtains the following.

\begin{proposition}
\label{prop:metrics-model-metric-density}
Let \(K\) be discretely valued and let \(X\) be an integral projective
\(K\)-scheme with a line bundle \(L\). Then every continuous metric on
\(L^{\mathrm{an}}\) is a uniform limit of model metrics; that is, model
metrics are \(d_\infty\)-dense in \(\mathcal M(L)\).
\end{proposition}

\begin{proof}
By \Cref{prop:line-bundle-extension} there are \(m\ge1\), a projective model
\(\mathcal X\), and a line bundle \(\mathcal L\) on \(\mathcal X\) with
\(\mathcal L|_X\simeq mL\); the pair \((\mathcal X,m^{-1}\mathcal L)\)
defines a model metric \(\|\cdot\|_0\). Let \(\|\cdot\|\) be a continuous
metric and \(\phi\) its relative weight, a continuous function by
\Cref{prop:metrics-torsor}. By
\Cref{prop:metrics-model-function-density} there are
\(\varphi_j\in\mathcal D(X)_{\mathbf Q}\) with
\(\|\varphi_j-\phi\|_\infty\to0\). Each \(\|\cdot\|_0\rme^{-\varphi_j}\) is
a model metric, by \Cref{prop:metric-difference-model} read in the reverse
direction: twisting a model metric by the model function of a vertical
\(\mathbf Q\)-Cartier divisor amounts to replacing \(\mathcal L\) by
\(\mathcal L\otimes\mathcal O(mD)\) on a model determining \(D\). By
\eqref{eq:metrics-dinfty-sup} these model metrics converge to
\(\|\cdot\|\).
\end{proof}

Third---and this is the warning that governs the rest of the
book---\emph{density says nothing about positivity}. The approximating model
functions produced above are maxima and minima of Fubini--Study differences;
nothing in the construction makes the corresponding model line bundles
relatively nef, and indeed a continuous metric that is not semipositive is
approximated by model metrics that cannot all be semipositive. The
semipositive theory therefore needs its own approximation theorems, which
are the subject of \cref{part:bfj}; the same warning is repeated after
\Cref{thm:gm-metric-comparison} and again in
\cref{sec:metrics-positivity-interface}.

\section{Semipositive metrics}
\label{sec:metrics-semipositive}

We now come to the positivity notion on which the whole book turns. In
complex geometry, a smooth metric is semipositive when its curvature form is
a positive \((1,1)\)-form, a pointwise condition. Here no forms are
available, and the substitute is a global one: a model metric is
semipositive when the model line bundle representing it has nonnegative
degree on every curve in the special fiber. The reader should think of this
as the exact analogue of the characterization of nef line bundles by degrees
on curves, applied fiberwise to the degeneration rather than to \(X\)
itself. The following facts, all proved below, justify the definition:
it is independent of the representing model; it is stable under the
operations that the theory needs; and on a curve it is literally the
statement that the associated measure of \cref{sec:metrics-cl-measure} is
positive.

Return first to a proper model \(\mathcal X\) over a discrete valuation ring.

\begin{definition}[Model semipositivity]
A line bundle \(\mathcal L\) on \(\mathcal X\) is
\emph{relatively nef} if
\[
                         \deg_{\mathcal L}(C)\geq0
\]
for every proper integral curve \(C\) contained in the special fiber. A
model metric is \emph{semipositive} if it is represented by a relatively
nef rational model line bundle.
\end{definition}

Pullback preserves relative nefness: a vertical curve either maps to a
point, in which case the pulled-back degree is zero, or maps onto a curve,
in which case the projection formula gives a nonnegative multiple of the
original degree. Conversely, the model test shows that semipositivity is
independent of the determining model. Both halves of this assertion are
proved in the next lemma, because a definition that depended on the chosen
model would be useless.

\begin{lemma}
\label{lem:metrics-nef-independence}
Let \(K\) be discretely valued and let \(X\) be an integral proper
\(K\)-scheme.
\begin{enumerate}
\item Let \(\rho:\mathcal X'\to\mathcal X\) be a proper surjective morphism
of normal proper models inducing the identity on \(X\), and let
\(\mathcal L\) be a line bundle on \(\mathcal X\). Then \(\mathcal L\) is
relatively nef if and only if \(\rho^*\mathcal L\) is.
\item If two pairs \((\mathcal X_1,m_1^{-1}\mathcal L_1)\) and
\((\mathcal X_2,m_2^{-1}\mathcal L_2)\), with normal models
\(\mathcal X_i\), induce the same metric on \(L^{\mathrm{an}}\), then
\(\mathcal L_1\) is relatively nef if and only if \(\mathcal L_2\) is.
\end{enumerate}
Consequently the semipositivity of a model metric depends only on the
metric.
\end{lemma}

\begin{proof}
(1) The implication from \(\mathcal L\) to \(\rho^*\mathcal L\) is
\Cref{prop:nef-pullback}. Conversely, assume \(\rho^*\mathcal L\)
relatively nef and let \(C\subset\mathcal X_s\) be an integral curve. By
Chow's lemma we may replace \(\mathcal X'\) by a dominating projective
model, which preserves the hypothesis by \Cref{prop:nef-pullback}. Since
\(\rho\) is proper and surjective, \(\rho^{-1}(C)\) is a closed subset of
\(\mathcal X'_s\) mapping onto \(C\), so it contains an integral closed
subvariety \(W\) mapping onto \(C\) with \(\dim W=1\): take an irreducible
component of \(\rho^{-1}(C)\) dominating \(C\) and cut it repeatedly with
hyperplane sections, which does not destroy surjectivity onto \(C\) because
every fiber of a surjection from an irreducible projective variety of
dimension \(d\ge2\) onto a curve has dimension \(\ge d-1\ge1\). Then
\(W\to C\) is a finite
surjective morphism of integral proper curves over \(\widetilde K\), of some
degree \(e\ge1\), and the projection formula
(\Cref{thm:model-projection-formula} in the form of
\cite[Proposition~2.5(c)]{Ful}) gives
\[
 0\le(\rho^*\mathcal L\cdot W)=e\,(\mathcal L\cdot C),
\]
whence \((\mathcal L\cdot C)\ge0\).

(2) Choose a normal proper model \(\mathcal X\) dominating both, by
\Cref{prop:common-domination} followed by normalization, and pull back both
line bundles; by (1) this changes neither nefness hypothesis. The bundles
\(\mathcal L_1^{\otimes m_2}\) and \(\mathcal L_2^{\otimes m_1}\) then both
extend \(m_1m_2L\) and induce the same metric, so by
\Cref{prop:metric-difference-model} they differ by
\(\mathcal O_{\mathcal X}(D)\) for a vertical Cartier divisor \(D\) with
\(\varphi_D\equiv0\); by
\Cref{prop:metrics-divisor-model-function}\,(2), \(D=0\). Hence
\(\mathcal L_1^{\otimes m_2}\simeq\mathcal L_2^{\otimes m_1}\), and since
degrees are multiplicative under tensor powers, one is relatively nef if and
only if the other is, if and only if \(\mathcal L_1\) is, if and only if
\(\mathcal L_2\) is.
\end{proof}

\begin{definition}[Continuous semipositivity]
A continuous metric on \(L\) is \emph{semipositive in the sense of Zhang} if
it is a uniform limit of semipositive model metrics.
\end{definition}

This is a global definition. It precedes, and must not be conflated with, the
local differential definitions of \cref{part:cld} or the singular psh metrics
of \cref{part:bfj}.

\begin{proposition}\label{prop:semipositive-stability}
Semipositive model metrics, and hence continuous semipositive metrics, have
the following properties.
\begin{enumerate}
\item Tensor products and positive rational powers are semipositive.
\item Pullback by a morphism of proper varieties is semipositive.
\item If two semipositive metrics on the same line bundle have relative
weights \(\phi_1,\phi_2\), then the metric with relative weight
\(\max\{\phi_1,\phi_2\}\) is semipositive.
\item Uniform limits of continuous semipositive metrics are semipositive.
\end{enumerate}
\par
\end{proposition}

\begin{proof}
The first two assertions follow from additivity of degrees and the projection
formula. The last assertion is built into the definition by a diagonal
approximation.

In detail, for (1) let \((\mathcal X_i,m_i^{-1}\mathcal L_i)\), \(i=1,2\),
be semipositive model metrics on \(L_1,L_2\). Choose a normal model
\(\mathcal X\) dominating both and pull the bundles back; by
\Cref{lem:metrics-nef-independence}\,(1) they remain relatively nef, and by
\Cref{prop:metrics-model-continuous}\,(3) they induce the same metrics. The
bundle \(\mathcal L_1^{\otimes m_2}\otimes\mathcal L_2^{\otimes m_1}\)
represents the tensor product metric on \(m_1m_2(L_1+L_2)\), and for every
integral curve \(C\subset\mathcal X_s\),
\[
 \bigl(\mathcal L_1^{\otimes m_2}\otimes\mathcal L_2^{\otimes m_1}\cdot C\bigr)
 =m_2(\mathcal L_1\cdot C)+m_1(\mathcal L_2\cdot C)\ge0 .
\]
A positive rational power \(r=a/b\) is handled by replacing
\((\mathcal X,m^{-1}\mathcal L)\) by \((\mathcal X,(bm)^{-1}
\mathcal L^{\otimes a})\), whose degrees are those of \(\mathcal L\)
multiplied by \(a>0\).

For (2), let \(f:Y\to X\) be a morphism of proper \(K\)-varieties and let
\((\mathcal X,m^{-1}\mathcal L)\) be a semipositive model metric on \(L\).
Choose any proper model \(\mathcal Y_0\) of \(Y\) and let \(\mathcal Y\) be
the flat closure of the graph of \(f\) in
\(\mathcal Y_0\times_R\mathcal X\), a proper model of \(Y\) carrying a
morphism \(F:\mathcal Y\to\mathcal X\) extending \(f\); that the flat
closure of the graph is again a model is
\Cref{lem:flat-closure-generic-fibre}. Then
\(F^*\mathcal L\) extends \(m\,f^*L\) and induces \(f^*\|\cdot\|\) by
functoriality of centers, as in
\Cref{prop:metrics-model-continuous}\,(3). For an integral curve
\(C\subset\mathcal Y_s\), either \(F(C)\) is a point, and
\((F^*\mathcal L\cdot C)=0\) because \(F^*\mathcal L|_C\) is then trivial,
or \(F(C)\) is an integral curve in \(\mathcal X_s\), and the projection
formula gives
\((F^*\mathcal L\cdot C)=\deg(C/F(C))\,(\mathcal L\cdot F(C))\ge0\). Hence
\(F^*\mathcal L\) is relatively nef.

For the maximum, note first that in multiplicative language it is the
pointwise minimum of the two norms. For model metrics, the minimum is again
a model metric after blowing up the sum of the two corresponding fractional
ideals; this is \Cref{lem:metrics-model-lattice}, applied to the relative
weights with respect to a common model metric. Its local semipositivity is
the maximum theorem for semipositive PL metrics. This nontrivial model test
is proved over arbitrary complete
nontrivially valued fields by Gubler--Martin; in the original discrete,
residue-characteristic-zero scope it is also
\cite[Section~5]{BFJ16}. Passing to uniform limits proves the continuous
case. A complete elementary proof in the case of curves, where the
statement becomes a convexity property of piecewise affine functions on a
metric graph, is given in \Cref{lem:metrics-max-graph} below.

For (4), let \(\|\cdot\|_j\) be continuous semipositive metrics converging
uniformly to \(\|\cdot\|\). For each \(j\), the definition provides a
semipositive model metric \(\|\cdot\|_j'\) with
\(d_\infty(\|\cdot\|_j,\|\cdot\|_j')\le1/j\); then
\[
 d_\infty(\|\cdot\|_j',\|\cdot\|)\le
 d_\infty(\|\cdot\|_j',\|\cdot\|_j)+d_\infty(\|\cdot\|_j,\|\cdot\|)
 \longrightarrow0 ,
\]
so \(\|\cdot\|\) is itself a uniform limit of semipositive model metrics.
\end{proof}

The proof of assertion (3) is the only one that is not elementary, and it is
worth isolating what is at stake. That the maximum of two model metrics is
again a model metric is \Cref{lem:metrics-model-lattice}; the difficulty is
that the model line bundle produced by the blowup need not obviously be
relatively nef, since the blowup introduces new components of the special
fiber and the degrees there have to be controlled. On curves this control is
a two-line convexity computation, which we now give in full because it is
also the source of the geometric picture in higher dimensions.

\begin{lemma}
\label{lem:metrics-max-graph}
Let \(\Gamma\) be a compact metric graph, let \(\mu_0\) be a positive Radon
measure on \(\Gamma\), and let \(\varphi_1,\varphi_2\) be piecewise affine
functions on \(\Gamma\) with
\[
 \mu_0+\Delta_\Gamma\varphi_i\ge0\qquad(i=1,2).
\]
Then \(\varphi:=\max\{\varphi_1,\varphi_2\}\) is piecewise affine and
\[
 \mu_0+\Delta_\Gamma\varphi\ge0 .
\]
\end{lemma}

\begin{proof}
That \(\varphi\) is piecewise affine is clear: subdivide \(\Gamma\) so that
both \(\varphi_i\) are affine on each edge and the finitely many points
where \(\varphi_1-\varphi_2\) changes sign are vertices. Since
\(\Delta_\Gamma\varphi\) is a finite atomic measure by
\eqref{eq:graph-laplacian}, it suffices to check, for each point
\(p\in\Gamma\), that
\[
 \mu_0(\{p\})+\Delta_\Gamma\varphi(\{p\})\ge0 .
\]
Fix \(p\) and choose \(i\) with \(\varphi_i(p)=\varphi(p)\); such an \(i\)
exists by definition of the maximum. For every tangent direction
\(\nu\in T_p\Gamma\), the inequality \(\varphi\ge\varphi_i\) near \(p\),
together with the equality at \(p\), gives
\[
 d_\nu\varphi(p)
 =\lim_{\varepsilon\downarrow0}
   \frac{\varphi(p+\varepsilon\nu)-\varphi(p)}{\varepsilon}
 \ge
 \lim_{\varepsilon\downarrow0}
   \frac{\varphi_i(p+\varepsilon\nu)-\varphi_i(p)}{\varepsilon}
 =d_\nu\varphi_i(p).
\]
Summing over the finitely many directions at \(p\) yields
\(\Delta_\Gamma\varphi(\{p\})\ge\Delta_\Gamma\varphi_i(\{p\})\), and
therefore
\[
 \mu_0(\{p\})+\Delta_\Gamma\varphi(\{p\})
 \ge\mu_0(\{p\})+\Delta_\Gamma\varphi_i(\{p\})\ge0 .
\]
\end{proof}

\begin{example}[A maximum of two model metrics on a skeleton]
\label{ex:metrics-max-skeleton}
\index{maximum!of model metrics}\index{semipositivity!stability under maximum}
Let \(\mathcal C=\{XY=\pi Z^2\}\subset\mathbf P^2_R\) be the conic model
\eqref{eq:running-conic}, with generic fiber the smooth conic \(C\),
\(L=\mathcal O(1)|_C\) of degree \(2\), and skeleton
\(\Gamma=[0,1]\) parametrized by \(s=v(x)\). The reference model metric has
curvature \(\mu_0=\delta_{v_0}+\delta_{v_1}\) by
\Cref{prop:running-conic-measure}.

Consider the two affine model functions
\[
 \varphi_1(s)=s-\tfrac12,\qquad \varphi_2(s)=\tfrac12-s .
\]
By \eqref{eq:running-perturbed-measure}, with \((A,B)=(-\tfrac12,\tfrac12)\)
and \((\tfrac12,-\tfrac12)\) respectively,
\[
 \mu_0+\Delta\varphi_1=2\delta_{v_0},
 \qquad
 \mu_0+\Delta\varphi_2=2\delta_{v_1},
\]
both positive, so both twisted metrics are semipositive; each lies on the
boundary of the nef cone, since \(|A-B|=1\) is the extreme case of
\eqref{eq:running-nef-interval}.

Now take the maximum, \(\varphi=\max\{\varphi_1,\varphi_2\}=|s-\tfrac12|\).
Its outgoing slopes are \(-1\) at \(s=0\), \(+1+1=+2\) at \(s=\tfrac12\),
and \(-1\) at \(s=1\), so
\[
 \Delta\varphi=-\delta_{v_0}+2\delta_{v_{1/2}}-\delta_{v_1},
\]
of total mass \(0\), and
\[
 \mu_0+\Delta\varphi=2\,\delta_{v_{1/2}}\ \ge0 ,
\]
of total mass \(2=\deg_CL\). The maximum is semipositive, as
\Cref{lem:metrics-max-graph} predicts, and the mechanism is visible: taking
a maximum creates a convex corner, whose Laplacian is a positive atom, and
positive atoms can only help.

\emph{The same computation at model level.} Blow up the node of
\(\mathcal C_s\), that is, the ideal \((x,y)\) at the point \(x=y=0\) of the
special fiber; the local equation \(xy=\pi\) shows that \(\mathcal C\) is
regular there, so the blowup \(\rho:\mathcal C'\to\mathcal C\) is a regular
model with special fiber
\[
 \mathcal C'_s=\widetilde C_0+2E+\widetilde C_1,
\]
the multiplicity \(2\) arising from
\(\operatorname{ord}_E(\pi)=\operatorname{ord}_E(x)+\operatorname{ord}_E(y)=2\).
The divisorial point of \(E\) is the midpoint \(v_{1/2}\), consistent with
the normalization \(x_E=b_E^{-1}\operatorname{ord}_E\). From
\(\mathcal C'_s\cdot C=0\) for every fibral curve \(C\) and the incidences
\((\widetilde C_0\cdot E)=(E\cdot\widetilde C_1)=1\),
\((\widetilde C_0\cdot\widetilde C_1)=0\), one computes the intersection
matrix
\[
 (\widetilde C_0^2)=(\widetilde C_1^2)=-2,\qquad (E^2)=-1 .
\]
The function \(\varphi=|s-\tfrac12|\) takes the values
\(\tfrac12,0,\tfrac12\) at the three divisorial points, so it is the model
function of the vertical \(\mathbf Q\)-divisor
\(D=\tfrac12\widetilde C_0+\tfrac12\widetilde C_1\), by
\Cref{prop:metrics-divisor-model-function}\,(2) together with
\(b_E=2\). With \(\mathcal L'=\rho^*\mathcal L\otimes\mathcal O(D)\) and
\(\deg(\rho^*\mathcal L|_{\widetilde C_i})=1\),
\(\deg(\rho^*\mathcal L|_E)=0\), we get
\[
 \begin{aligned}
 \deg(\mathcal L'|_{\widetilde C_0})
 &=1+\tfrac12(-2)+0=0,\\
 \deg(\mathcal L'|_{E})&=0+\tfrac12\cdot1+\tfrac12\cdot1=1,\\
 \deg(\mathcal L'|_{\widetilde C_1})&=1+0+\tfrac12(-2)=0 .
 \end{aligned}
\]
The Chambert--Loir measure of \cref{sec:metrics-cl-measure} is therefore
\[
 1\cdot0\cdot\delta_{v_0}+2\cdot1\cdot\delta_{v_{1/2}}
 +1\cdot0\cdot\delta_{v_1}
 =2\,\delta_{v_{1/2}},
\]
in exact agreement with the Laplacian computation, multiplicity \(b_E=2\)
included. All degrees are nonnegative, so \(\mathcal L'\) is relatively nef
and the maximum metric is semipositive on the nose---no approximation and
no citation are needed in this case.
\end{example}

\begin{theorem}\label{thm:model-semipositive-closed}
Let \(X\) be proper over a complete nontrivially valued field. If a model
metric is a pointwise limit, on all divisorial points, of semipositive model
metrics, then it is itself a semipositive model metric.
\end{theorem}
\begin{proof}[Foundational closedness input]
This is the closedness theorem of Gubler--Martin
\cite[Theorems~1.3 and 5.5]{GM19}; the ample base-ideal step in its proof is
\cite[Proposition~5.1]{GM19}.  Their criterion tests the pointwise limit on
divisorial points and concludes that the limiting model metric is represented
by a relatively nef model line bundle.  These are exactly the hypothesis and
conclusion above, so the cited theorem applies directly.

\end{proof}

\section{Three positivity languages}
\label{sec:metrics-positivity-interface}

Three notions meet at the model level.

\begin{enumerate}
\item \emph{Global model positivity}: a model line bundle is relatively nef.
This is the Zhang--BFJ language used in this chapter.
\item \emph{Delta-current positivity}: Gubler--Künnemann attach a first Chern
delta-current to a formal metric and compare its positivity with relative
nefness.
\item \emph{Curve-test positivity}: Gubler--Rabinoff study semipositive
piecewise linear functions and classically psh functions using pullback to
analytic curves.
\end{enumerate}

The three are not competing definitions of one notion; they are conditions
of different logical shapes, and it is useful to see the shapes side by
side.

\medskip
\noindent
\begin{tabular}{@{}p{.24\textwidth}p{.34\textwidth}p{.34\textwidth}@{}}
\toprule
language & datum tested & shape of the condition\\
\midrule
model positivity
 & a model line bundle \(\mathcal L\)
 & finitely many degree inequalities on curves in \(\mathcal X_s\)\\
delta-current positivity
 & a formal metric and its first Chern delta-current
 & positivity of a current on tropical charts\\
curve-test positivity
 & a piecewise linear function or metric
 & subharmonicity after pullback to analytic curves\\
\bottomrule
\end{tabular}
\medskip

\noindent The first is global and combinatorial, the second is local and
differential, and the third is local but tested along one-dimensional
subvarieties. Each is stable under the operations of
\Cref{prop:semipositive-stability} for its own reasons, and none of the
three is defined for arbitrary singular potentials without further work.

On proper algebraic spaces and for formal or model metrics, the published
comparison theorems identify the appropriate positivity conditions. Their
hypotheses will be stated in
\crefrange{chap:cld-metrized-bundles}{chap:cld-examples}. We do not infer from them an
unrestricted equivalence among arbitrary CLD-psh functions, arbitrary
BFJ singular metrics, and arbitrary curve-test psh functions. In particular,
Gubler--Martin's density of PL metrics does not by itself preserve
semipositivity, and a local comparison at model metrics does not identify
the completions used by different theories.

\section{Chambert--Loir measures of model metrics}
\label{sec:metrics-cl-measure}

We can now attach a measure to model metrics. The construction is forced by
two requirements. It must be multilinear in the metrized line bundles, like
an intersection product; and its total mass must be the intersection number
\((L_1\cdots L_n)\), so that it deserves to be called a curvature. On a
model, the only available multilinear data are the degrees of the model line
bundles on the components of the special fiber, and the only available
points are the divisorial ones. The definition below is the unique way to
combine them consistently, and the multiplicity \(b_Y\) is not a decoration:
it is exactly what makes the total mass come out right, as the proof of
\Cref{thm:cl-model-properties} shows.

Assume for this section that \(K\) is discretely valued. Let \(X\) be
projective of pure dimension \(n\), and let \(\mathcal X\) be a normal
projective \(R\)-model. Write
\[
                        \mathcal X_s=\sum_Y b_YY,
\]
where \(Y\) ranges over the irreducible components. Let \(x_Y\) be the
divisorial point defined by \(b_Y^{-1}\operatorname{ord}_Y\).

For \(1\leq i\leq n\), let \(L_i\) be a line bundle on \(X\), and let
\(\|\cdot\|_i\) be a model metric determined by a line bundle
\(\mathcal L_i\) extending \(L_i^{\otimes m_i}\), with \(m_i>0\).

\begin{definition}[Mixed model measure]
The \emph{Chambert--Loir measure} of these model metrics is
\[
\begin{split}
 &c_1(L_1,\|\cdot\|_1)\wedge\cdots\wedge
                 c_1(L_n,\|\cdot\|_n)\\
 &\quad =
 \frac{1}{m_1\cdots m_n}
 \sum_Y b_Y
 \left(
 c_1(\mathcal L_1|_Y)\cdots c_1(\mathcal L_n|_Y)
 \right)\delta_{x_Y}.
\end{split}
\]
The intersection number in parentheses is computed on the
\(n\)-dimensional proper \(\widetilde K\)-scheme \(Y\).
\end{definition}

The multiplicity \(b_Y\) is part of the formula. It compensates for the
normalization \(v_{x_Y}=b_Y^{-1}\operatorname{ord}_Y\) and is also forced by
the total-mass calculation below.

\begin{theorem}\label{thm:cl-model-properties}
The mixed model measure is independent of the common normal model and of the
chosen tensor powers. It is symmetric and multilinear in the metrized line
bundles. Its total mass is
\[
 \int_{X^{\mathrm{an}}}
 c_1(L_1,\|\cdot\|_1)\wedge\cdots\wedge
 c_1(L_n,\|\cdot\|_n)
                  =(L_1\cdots L_n).
\]
If every model metric is semipositive, the measure is positive.
\end{theorem}

\begin{proof}
Suppose first that \(p:\mathcal X'\to\mathcal X\) is a dominating normal
model and replace every \(\mathcal L_i\) by \(p^*\mathcal L_i\). A component
of \(\mathcal X'_s\) that maps into a subset of dimension less than \(n\)
has zero intersection against \(n\) pulled-back line bundles. A component
dominating a component \(Y\) contributes the expected degree by the
projection formula. Including special-fiber multiplicities, the sum of the
contributions over \(Y\) is unchanged. Any two models admit a common
dominating model, proving independence. Replacing \(\mathcal L_i\) by a
tensor power multiplies the corresponding intersection number and the
denominator by the same factor.

The first paragraph deserves one more line of detail, since it is where the
multiplicities are checked.  Under the birational domination
\(p:\mathcal X'\to\mathcal X\), a component \(Y\subset\mathcal X_s\) has a
unique strict transform \(Y'\) dominating it.  The map \(Y'\to Y\) is
birational, \(b_{Y'}=b_Y\), and \(x_{Y'}=x_Y\), since the two normalized
orders define the same valuation of \(K(X)\).  Every other component of
\(\mathcal X'_s\) maps to a subset of dimension less than \(n\), so its
intersection against the \(n\) pulled-back line bundles vanishes by the
projection formula.  Thus \Cref{thm:model-projection-formula} gives equality
of the two atomic measures, atom by atom.

Symmetry and multilinearity are inherited from intersection theory. For the
total mass, use the equality of cycles
\(\mathcal X_s=\sum_Yb_YY\):
\[
\begin{split}
 \sum_Yb_Y
  (c_1(\mathcal L_1)\cdots c_1(\mathcal L_n)|Y)
 &=(c_1(\mathcal L_1)\cdots c_1(\mathcal L_n)|\mathcal X_s)\\
 &=m_1\cdots m_n(L_1\cdots L_n).
\end{split}
\]
The last equality is constancy of intersection degree in the flat family.
Finally, mixed intersections of nef line bundles on each component \(Y\)
are nonnegative, so all atomic coefficients are nonnegative.
\end{proof}

\begin{remark}
For arbitrary model metrics, the formula defines a signed measure. The
symbol \(c_1\) here records multilinearity and compatibility with
intersection theory. \Cref{chap:cld-metrized-bundles} constructs actual CLD first Chern currents
and proves that their wedge product recovers this measure on the common
formal-model domain.
\end{remark}

The next statement is the computational heart of the whole model theory: it
identifies the integral of a model function against a model measure with a
vertical intersection number. Everything later---the extension to continuous
metrics, the energy functionals of \cref{part:bfj}, the differentiability
statements of \cref{chap:bfj-differentiability}---is obtained by
manipulating this one identity.

\begin{proposition}
\label{prop:metrics-cl-model-function-integral}
Let \(K\) be discretely valued, let \(X\) be projective of pure dimension
\(n\), and let \(\mathcal X\) be a normal projective model carrying line
bundles \(\mathcal L_1,\ldots,\mathcal L_n\) with
\(\mathcal L_i|_X\simeq L_i^{\otimes m_i}\). Let \(D\) be a vertical
\(\mathbf Q\)-Cartier divisor on \(\mathcal X\), with model function
\(\varphi_D\). Then
\begin{equation}
\label{eq:metrics-cl-pairing}
 \int_{X^{\mathrm{an}}}\varphi_D\;
 c_1(L_1,\|\cdot\|_1)\wedge\cdots\wedge c_1(L_n,\|\cdot\|_n)
 =\frac{(D\cdot\mathcal L_1\cdots\mathcal L_n)_{\mathcal X}}
        {m_1\cdots m_n}.
\end{equation}
In particular, taking \(D=\mathcal X_s\) recovers the total mass
\((L_1\cdots L_n)\), and the pairing on the left depends on \(D\) only
through \(\varphi_D\).
\end{proposition}

\begin{proof}
Write \(\mathcal X_s=\sum_Yb_YY\) and \(D=\sum_Ya_YY\). The measure is
atomic, supported on the finitely many divisorial points \(x_Y\), so
\[
 \int\varphi_D\,c_1(L_1,\|\cdot\|_1)\wedge\cdots
 =\frac1{m_1\cdots m_n}\sum_Yb_Y\,\varphi_D(x_Y)\,
   \bigl(\mathcal L_1|_Y\cdots\mathcal L_n|_Y\bigr).
\]
By \Cref{prop:metrics-divisor-model-function}\,(2) we have
\(\varphi_D(x_Y)=a_Y/b_Y\), so the multiplicities cancel and the sum becomes
\(\sum_Ya_Y(\mathcal L_1|_Y\cdots\mathcal L_n|_Y)\), which is
\((D\cdot\mathcal L_1\cdots\mathcal L_n)_{\mathcal X}\) by
\Cref{lem:vertical-intersection-degree}. For \(D=\mathcal X_s\) one has
\(\varphi_D\equiv1\) and the right-hand side is \((L_1\cdots L_n)\) by
\Cref{cor:special-fibre-relation}. The last assertion follows because the
left-hand side manifestly depends only on the function \(\varphi_D\), while
the right-hand side is an intersection number; combined with
\Cref{prop:metrics-divisor-model-function}\,(2) this also shows that the two
sides are unchanged under passage to a dominating model, as
\Cref{thm:model-projection-formula} requires.
\end{proof}

The cancellation of \(b_Y\) in this proof is the precise sense in which the
multiplicity in the definition of the measure and the normalization
\(x_Y=b_Y^{-1}\operatorname{ord}_Y\) of the divisorial point are two halves
of one convention. Using \(\operatorname{ord}_Y\) with the same coefficients,
or \(x_Y\) without the coefficient \(b_Y\), would break
\eqref{eq:metrics-cl-pairing}, and with it every energy formula in the book.
This is the content of the warning after \eqref{eq:appJ-cl-measure}.

\begin{example}[Good reduction]
Let \(X=\mathbf P^n_K\),
\(\mathcal X=\mathbf P^n_R\), and give \(L=\mathcal O(1)\) its standard
model metric. The special fiber is irreducible and reduced, and
\((c_1(\mathcal O(1))^n|\mathbf P^n_{\widetilde K})=1\). Hence
\[
                   c_1(\mathcal O(1),\|\cdot\|_{\mathrm{std}})^n
                   =\delta_{x_{\mathcal X_s}}.
\]
For \(n=1\), this is the Dirac mass at the Gauss point of the projective
line.
\end{example}

\begin{example}[A semistable curve]
Let \(X\) be a smooth projective curve and let \(\mathcal X\) be a regular
semistable model, so every \(b_Y=1\). For a model line bundle
\(\mathcal L\),
\[
                 c_1(L,\|\cdot\|_{\mathcal L})
                    =\sum_Y\deg(\mathcal L|_Y)\,\delta_{x_Y}.
\]
Semipositivity is exactly the nonnegativity of these degrees, and their sum
is \(\deg L\). The redistribution of these atomic masses under refinement
is the intersection-theoretic ancestor of the graph Laplacian in
\cref{chap:model-calculations}.
\end{example}

\begin{example}[One blowup of the projective line, at model level]
\label{ex:metrics-cl-p1-blowup}
We return to \cref{ex:metrics-p1-blowup} and recompute its measure from
intersection numbers, as promised there. Let \(X=\mathbf P^1_K\),
\(L=\mathcal O(1)\), and let \(\rho:\mathcal X'\to\mathbf P^1_R\) be the
blowup of the closed point \(p\) of the special fiber given by \(T\equiv0\).
The model \(\mathcal X'\) is a regular projective surface over \(R\) with
\[
 \mathcal X'_s=\widetilde C+E,\qquad b_{\widetilde C}=b_E=1,
\]
and \((E^2)=-1\), \((\widetilde C\cdot E)=1\); the value
\((\widetilde C^2)=-1\) then follows from
\(\mathcal X'_s\cdot\widetilde C=0\). Take
\(\mathcal L'=\rho^*\mathcal O(1)\otimes\mathcal O_{\mathcal X'}(-E)\), whose
generic fiber is \(\mathcal O(1)\) and whose metric was computed in
\cref{ex:metrics-p1-blowup} to have weight \(\max\{-1,\log|T|\}\). Then
\[
 \deg\bigl(\mathcal L'|_{E}\bigr)=0-(E^2)=1,
 \qquad
 \deg\bigl(\mathcal L'|_{\widetilde C}\bigr)
 =1-(E\cdot\widetilde C)=0,
\]
using that \(\rho^*\mathcal O(1)\) has degree \(0\) on the contracted curve
\(E\) and degree \(1\) on \(\widetilde C\). Hence
\[
 c_1\bigl(\mathcal O(1),\|\cdot\|_{\mathcal L'}\bigr)
 =1\cdot0\cdot\delta_{x_{\widetilde C}}+1\cdot1\cdot\delta_{x_E}
 =\delta_{x_E},
 \qquad x_E=\zeta_{0,|\pi|},
\]
of total mass \(1=\deg\mathcal O(1)\), in agreement with the Laplacian
computation \eqref{eq:metrics-p1-blowup-laplacian}. Both degrees are
nonnegative, so \(\mathcal L'\) is relatively nef and the metric is
semipositive; this is consistent with its description in
\cref{ex:metrics-fs-p1-written} as the Fubini--Study metric of the sections
\(\pi Z_0,Z_1\), which is semipositive by
\Cref{prop:metrics-fs-model}.
\end{example}

\begin{example}[A curve: the semistable chain, with all masses]
\label{ex:metrics-cl-curve-chain}
\index{Chambert--Loir measure!computed on a curve}
Let \(\mathcal C^{(2)}=\{XY=\pi^2Z^2\}\subset\mathbf P^2_R\) and let
\(\widetilde{\mathcal C}\to\mathcal C^{(2)}\) be the semistable chain model
of \Cref{ex:snc-chain-models}, with special fiber the chain
\(C_0,E,C_2\) of projective lines, all multiplicities one, whose divisorial
points sit at the three integral points \(0,1,2\) of the skeleton
\(\Gamma=[0,2]\). Let \(L=\mathcal O(1)|_C\), of degree \(2\) on the generic
fiber, and let
\(\mathcal L_0=\rho^*\bigl(\mathcal O_{\mathbf P^2_R}(1)|_{\mathcal
C^{(2)}}\bigr)\), whose fiber degrees are
\[
 \deg(\mathcal L_0|_{C_0})=1,\quad
 \deg(\mathcal L_0|_{E})=0,\quad
 \deg(\mathcal L_0|_{C_2})=1 ,
\]
as computed in \Cref{ex:degree-conservation-conic}. The associated measure
is therefore
\[
 \mu_0=\delta_{v_0}+\delta_{v_2},
\]
of mass \(2=\deg_CL\), the middle component carrying no mass because it is
contracted.

Now twist by a vertical divisor \(D=a_0C_0+a_1E+a_2C_2\), whose model
function takes the values \(a_0,a_1,a_2\) at the three divisorial points.
Using the intersection numbers computed in
\Cref{ex:chain-nef-functions},
\[
 (C_0^2)=(C_2^2)=-1,\quad (E^2)=-2,\quad
 (C_0\cdot E)=(E\cdot C_2)=1,\quad (C_0\cdot C_2)=0,
\]
the degrees of \(\mathcal L=\mathcal L_0\otimes\mathcal O(D)\) are
\[
 \begin{aligned}
 \deg(\mathcal L|_{C_0})&=1-a_0+a_1,\\
 \deg(\mathcal L|_{E})&=a_0-2a_1+a_2,\\
 \deg(\mathcal L|_{C_2})&=1+a_1-a_2,
 \end{aligned}
\]
so that
\begin{equation}
\label{eq:metrics-cl-chain-measure}
 c_1(L,\|\cdot\|_{\mathcal L})
 =(1+a_1-a_0)\delta_{v_0}
 +(a_0+a_2-2a_1)\delta_{v_1}
 +(1+a_1-a_2)\delta_{v_2}.
\end{equation}
The three masses sum to \(2=\deg_CL\) for every choice of \(a\), as
\Cref{thm:cl-model-properties} requires and as
\Cref{cor:special-fibre-relation} predicts.

\emph{Cross-check by the Laplacian.} Let \(\varphi\) be the piecewise affine
function on \(\Gamma=[0,2]\) with \(\varphi(i)=a_i\), affine on \([0,1]\) and
on \([1,2]\); it is the model function of \(D\). Summing outgoing slopes as
in \eqref{eq:graph-laplacian},
\[
 \Delta_\Gamma\varphi
 =(a_1-a_0)\delta_{v_0}
 +\bigl[(a_0-a_1)+(a_2-a_1)\bigr]\delta_{v_1}
 +(a_1-a_2)\delta_{v_2},
\]
and \(\mu_0+\Delta_\Gamma\varphi\) is exactly
\eqref{eq:metrics-cl-chain-measure}. Thus the graph Laplacian of
\cref{chap:model-calculations} and the intersection theory of
\cref{chap:models} agree on the nose, with the signs of
\cref{sec:appJ-model-functions}. Semipositivity holds precisely when the
three coefficients are nonnegative, that is, when \(\varphi\) is convex at
the interior vertex and its outgoing slopes at the two endpoints are at
least \(-1\); this is the nef cone computed in
\Cref{ex:chain-nef-functions}.
\end{example}

\begin{example}[A toric surface: the blowup of a fixed point of the special
fiber]
\label{ex:metrics-cl-toric-surface}
\index{Chambert--Loir measure!computed on a toric surface}
\index{toric geometry!model metric}
Let \(K\) be discretely valued with \(v(\pi)=1\), let
\(X=\mathbf P^2_K\) with homogeneous coordinates \(Z_0,Z_1,Z_2\), torus
coordinates \(x=Z_1/Z_0\), \(y=Z_2/Z_0\), and let \(L=\mathcal O(1)\), so
that \((L^2)=1\). The polytope of \(L\) is the standard simplex
\(P=\operatorname{conv}\{0,e_1,e_2\}\subset M_{\mathbf R}\), with
\(2!\operatorname{vol}_M(P)=1=(L^2)\).

Consider the Fubini--Study metric attached to the three sections
\[
 s_0=\pi Z_0,\qquad s_1=Z_1,\qquad s_2=Z_2
\]
of \(\mathcal O(1)\), with zero weights. They have no common zero, so
\Cref{prop:metrics-fs-model} applies: the metric is a semipositive model
metric. In the frame \(Z_0\), and in the tropical coordinates
\(u=(v(x),v(y))\) of \eqref{eq:toric-trop}, its weight is
\begin{equation}
\label{eq:metrics-toric-weight}
 F(u)=\max\{-1,\,-u_1,\,-u_2\},
\end{equation}
since \(\log|\pi|=-1\), \(\log|x|=-u_1\), and \(\log|y|=-u_2\). It is
convex, piecewise affine, and \(F-\check h_P\) is bounded, where
\(\check h_P(u)=\max\{0,-u_1,-u_2\}\) is the recession function of
\cref{sec:toric-convex}; indeed \(-1\le F-\check h_P\le0\).

\emph{The model.} Clearing denominators, the metric is induced by the ideal
\(\mathfrak a=(\pi,x,y)\) of the torus-fixed point \(p\) of the special
fiber \(\mathbf P^2_{\widetilde K}\), namely the reduction of
\([1:0:0]\). Let \(\rho:\mathcal X\to\mathbf P^2_R\) be the blowup of
\(\mathfrak a\), that is, of the closed point \(p\) of the three-fold
\(\mathbf P^2_R\). Its exceptional divisor is
\(E\simeq\mathbf P^2_{\widetilde K}\), with
\(\mathcal O(-E)|_E\simeq\mathcal O_{\mathbf P^2}(1)\), and the strict
transform \(\widetilde{X_s}\) of the special fiber is the blowup of
\(\mathbf P^2_{\widetilde K}\) at \(p\), a Hirzebruch surface. Since
\(\pi\in\mathfrak m_p\setminus\mathfrak m_p^2\), both components occur with
multiplicity one:
\[
 \mathcal X_s=\widetilde{X_s}+E .
\]
The model line bundle is
\(\mathcal L=\rho^*\mathcal O(1)\otimes\mathcal O_{\mathcal X}(-E)\), and
the corresponding model function of \(-E\) is
\(-\varphi_E=-\min\{1,u_1^+,u_2^+\}\), which added to \(\check h_P\) gives
\eqref{eq:metrics-toric-weight}.

\emph{The masses.} Write \(h\) for the pullback of the hyperplane class of
\(\mathbf P^2_{\widetilde K}\) to \(\widetilde{X_s}\) and \(e\) for the
exceptional curve of \(\widetilde{X_s}\to\mathbf P^2_{\widetilde K}\), so
that \((h^2)=1\), \((h\cdot e)=0\), \((e^2)=-1\). On the two components,
\[
 \mathcal L|_{E}\simeq\mathcal O_{\mathbf P^2}(1),
 \qquad
 \mathcal L|_{\widetilde{X_s}}\simeq h-e,
\]
the first because \(\rho^*\mathcal O(1)\) is trivial on the contracted
divisor \(E\) and \(\mathcal O(-E)|_E=\mathcal O_{\mathbf P^2}(1)\), the
second because \(\mathcal O(-E)|_{\widetilde{X_s}}=-e\). Hence
\[
 \bigl(\mathcal L|_{E}\bigr)^2=1,
 \qquad
 \bigl(\mathcal L|_{\widetilde{X_s}}\bigr)^2=(h-e)^2=1-1=0,
\]
and the Chambert--Loir measure is
\begin{equation}
\label{eq:metrics-toric-measure}
 c_1\bigl(\mathcal O(1),\|\cdot\|_{\mathcal L}\bigr)^2
 =1\cdot0\cdot\delta_{x_{\widetilde{X_s}}}+1\cdot1\cdot\delta_{x_E}
 =\delta_{x_E},
\end{equation}
of total mass \(1=(L^2)\), as \Cref{thm:cl-model-properties} requires. The
divisorial point \(x_E\) is the monomial valuation with
\(v(x)=v(y)=v(\pi)=1\), that is, the point \(u=(1,1)\) of the toric
skeleton.

\emph{Cross-check by convex geometry.} The function
\eqref{eq:metrics-toric-weight} is the maximum of the three affine functions
with slopes \(0,-e_1^*,-e_2^*\); their common domain of coincidence is the
single point \(u=(1,1)\), and there
\[
 \partial F(1,1)=\operatorname{conv}\{0,-e_1,-e_2\}=-P,
 \qquad
 \operatorname{vol}_M(-P)=\tfrac12 .
\]
By \eqref{eq:pl-real-ma}, \(\operatorname{MA}_{\mathbf R}(F)
=\tfrac12\delta_{(1,1)}\), and
\Cref{thm:toric-ma-comparison} gives
\[
 \operatorname{trop}_*c_1(L,\|\cdot\|)^2
 =2!\operatorname{MA}_{\mathbf R}(F)
 =\delta_{(1,1)},
\]
which is \eqref{eq:metrics-toric-measure} transported to
\(N_{\mathbf R}\). The two computations---one by intersection numbers on the
two components of a special fiber, one by the volume of a subgradient
polytope---agree, including the factor \(n!=2\) of
\Cref{lem:toric-real-total-mass}, and both give the total mass
\((L^2)=1\).

\emph{Nefness.} On \(E\simeq\mathbf P^2\) the bundle \(\mathcal O(1)\) is
ample. On the Hirzebruch surface \(\widetilde{X_s}\), the class \(h-e\) is
nef: it has degree \((h-e)\cdot e=1\ge0\) on the exceptional curve,
degree \((h-e)^2=0\ge0\) on the strict transform of a line through \(p\),
and degree \((h-e)\cdot h=1\ge0\) on a general line, and these classes
generate the cone of curves. Hence \(\mathcal L\) is relatively nef and the
metric is semipositive, as it must be, being a Fubini--Study metric. The
example also shows how mass migrates: the standard metric puts its unit mass
at the Gauss point \(u=(0,0)\), whereas this one puts it at \(u=(1,1)\).
\end{example}

\section{Continuous semipositive measures}
\label{sec:metrics-continuous-cl}

Continue to assume that \(K\) is discretely valued and that \(X\) is an
integral projective scheme of pure dimension \(n\).  Let
\(\|\cdot\|_i\) now be continuous semipositive metrics. Choose
semipositive model metrics \(\|\cdot\|_{i,j}\) converging uniformly to them.
The strategy of the following proof is to extend a multilinear operation
from a dense subspace.  The nontrivial input is the local-height continuity
estimate \eqref{eq:cl-model-cauchy-bound}, imported from Chambert--Loir.
This is why the extension theorem is stated for
\emph{semipositive} metrics and does not extend verbatim to arbitrary
continuous ones.

\begin{theorem}\label{thm:cl-continuous-extension}
The model measures
\[
 c_1(L_1,\|\cdot\|_{1,j})\wedge\cdots\wedge
 c_1(L_n,\|\cdot\|_{n,j})
\]
converge weakly to a positive Radon measure that is independent of all
approximating sequences. The resulting mixed measure is symmetric,
multilinear, has total mass \((L_1\cdots L_n)\), and is continuous under
uniform convergence of semipositive metrics.
\end{theorem}
\begin{proof}[Chambert--Loir local-height input]
We use the convergence and independence theorem
\cite[Proposition~2.7]{CL06}, proved there from the local-height recursion of
\cite[\S\S2.3--2.4]{CL06}.  In the present notation, the uniform Cauchy
estimate in that proof is the
following.  For every rational model function \(f\), there is a constant
\(C_f\), depending only on \(f\) and the numerical classes \(L_i\), such
that two semipositive model-metric tuples \(\boldsymbol\eta\) and
\(\boldsymbol\eta'\) satisfy
\begin{equation}\label{eq:cl-model-cauchy-bound}
 \left|\int f\,\rd\mu_{\boldsymbol\eta}
       -\int f\,\rd\mu_{\boldsymbol\eta'}\right|
 \le C_f\sum_{i=1}^n d_\infty(\eta_i,\eta_i').
\end{equation}
This estimate is the imported step.  Its proof uses the local-height
recursion and integration by parts; the curvature of \(f\) can be signed,
so it is not replaced here by an integral against a purported positive
lower-degree measure.

Uniform convergence of the approximating metrics and
\eqref{eq:cl-model-cauchy-bound} show that the integrals against every model
function form a Cauchy sequence.  Rational model functions are uniformly
dense in the continuous functions on the compact analytification by
\Cref{prop:metrics-model-function-density}.  All model measures are positive and have
the fixed mass $(L_1\cdots L_n)$, so approximation of an arbitrary
$g\in C(X^{\mathrm{an}})$ by a model function extends convergence to $g$.
The Riesz representation theorem gives a positive Radon limit measure.

Apply the same estimate to two different approximating families, after
interleaving them.  Their integrals against every model function have the
same limit, hence the two measures coincide.  Symmetry and multilinearity
hold for model intersection numbers and pass to the limit.  Conservation of
degree on the special fibre gives the stated total mass at the model level,
so it also passes to the limit.  Finally, if the limiting continuous metrics
themselves vary uniformly, choose model approximants within a common uniform
error and use \eqref{eq:cl-model-cauchy-bound}; this proves weak continuity.

\end{proof}

For one semipositive metric on an ample line bundle \(L\), we abbreviate
\[
               \operatorname{MA}(\|\cdot\|)
               =c_1(L,\|\cdot\|)^n.
\]
Its total mass is \((L^n)\). We will explicitly divide by \((L^n)\) whenever
a probability measure is intended.

\begin{remark}[Where singular theory begins]
Uniform approximation only reaches continuous semipositive metrics. The BFJ
Monge--Amp\`ere operator for bounded and finite-energy singular metrics is a
new construction, not a formal consequence of
\Cref{thm:cl-continuous-extension}. Similarly, the local
Bedford--Taylor products of CLD require their own locality and approximation
theorems.
\end{remark}

\section{Operations and change of reference}
\label{sec:metrics-operations}

Let \(\|\cdot\|_0\) be a reference metric on \(L\), and write a second metric
as \(\|\cdot\|_0\rme^{-\phi}\). Then:
\[
\begin{array}{ccl}
\text{tensor product} &\longleftrightarrow& \text{addition of weights},\\
\text{dual metric} &\longleftrightarrow& \text{negation of weights},\\
\text{pullback} &\longleftrightarrow& \text{composition with the map},\\
\text{pointwise minimum of norms}
  &\longleftrightarrow& \text{maximum of weights}.
\end{array}
\]
Changing the reference metric from \(\|\cdot\|_0\) to
\(\|\cdot\|_0\rme^{-\psi}\) replaces the relative weight \(\phi\) by
\(\phi-\psi\); the underlying metric does not change. Every later use of the
terms potential, maximum, or model function is relative to this convention.

Because this bookkeeping is invisible in the notation and fatal when done
wrongly, we state it once as a proposition and use it silently afterwards.

\begin{proposition}
\label{prop:metrics-change-of-reference}
Let \(X\) be projective of pure dimension \(n\) over a discretely valued
\(K\), let \(L_1,\ldots,L_n\) be line bundles with continuous semipositive
metrics, and let \(\|\cdot\|_{0,i}\) be reference metrics on \(L_i\), with
relative weights \(\phi_i\). Replace \(\|\cdot\|_{0,i}\) by
\(\|\cdot\|_{0,i}\rme^{-\chi_i}\) for continuous \(\chi_i\).
\begin{enumerate}
\item The relative weights become \(\phi_i'=\phi_i-\chi_i\), and every
inequality, uniform estimate, or maximum formed from weights with respect to
a common reference is unchanged.
\item The mixed measure
\(c_1(L_1,\|\cdot\|_1)\wedge\cdots\wedge c_1(L_n,\|\cdot\|_n)\) is
unchanged: it is an invariant of the metrics, not of the weights. In
particular, writing \(\MA(\phi)\) for the measure attached to the metric of
relative weight \(\phi\) with respect to a fixed reference, one has
\(\MA'(\phi-\chi)=\MA(\phi)\).
\item Adding a constant \(c\in\mathbf R\) to one weight does not change the
mixed measure, and changes the pairing with a fixed continuous function
\(g\) by nothing at all; conversely, for a fixed measure of total mass
\((L_1\cdots L_n)\),
\[
 \int g'\,c_1(L_1,\|\cdot\|_1)\wedge\cdots
 =\int g\,c_1(L_1,\|\cdot\|_1)\wedge\cdots
   +c\,(L_1\cdots L_n)
 \qquad\text{if }g'=g+c .
\]
\item Semipositivity of \(\|\cdot\|_i\) is a property of the metric; the
weight \(\phi_i\) is semipositive relative to \(\|\cdot\|_{0,i}\) if and only
if \(\phi_i-\chi_i\) is semipositive relative to the new reference.
\end{enumerate}
\end{proposition}

\begin{proof}
(1) is \Cref{prop:metrics-torsor}\,(2). (2) holds because the measure was
defined in \cref{sec:metrics-cl-measure} from model line bundles
representing the metrics, with no reference metric entering; the displayed
identity is a restatement of (1). (3) The first assertion follows from (2),
since adding a constant to a weight changes the metric by an overall
positive scalar, which is represented on a model by adding a multiple of
\(\mathcal X_s\), a divisor whose degrees on the components of the special
fiber vanish by \Cref{cor:special-fibre-relation} applied to
\(\mathcal X_s=\operatorname{div}(\pi)\); the second is linearity of the
integral together with the total-mass formula of
\Cref{thm:cl-model-properties}. (4) is
\Cref{lem:metrics-nef-independence} together with the observation that the
notion of semipositivity refers to the metric \(\|\cdot\|_{0,i}\rme^{-\phi_i}
=\bigl(\|\cdot\|_{0,i}\rme^{-\chi_i}\bigr)\rme^{-(\phi_i-\chi_i)}\), which is
the same metric written twice.
\end{proof}

Three consequences are used constantly. Statements about a single potential,
such as ``\(\phi\) is semipositive'' or ``\(\MA(\phi)=\mu\)'', are
statements about a metric and are reference-independent once the reference
is displayed. Statements comparing two potentials with respect to the same
reference, such as \(\phi_1\le\phi_2\) or
\(\|\phi_1-\phi_2\|_\infty\le\varepsilon\), are reference-independent
outright. But a numerical statement about a single potential, such as
\(\sup\phi=0\) or \(E(\phi)=0\), depends on the reference and must never be
quoted without it; this is why energies are normalized to vanish at a named
reference metric, as recorded in \cref{sec:appJ-energy}.

For a generically finite morphism \(f:Y\to X\) of degree \(d\), the model
projection formula gives
\[
 \begin{aligned}
 &f_*\!\left(
 c_1(f^*L_1,f^*\|\cdot\|_1)\wedge\cdots\wedge
 c_1(f^*L_n,f^*\|\cdot\|_n)\right)\\
 &\hspace{5em}=
 d\,c_1(L_1,\|\cdot\|_1)\wedge\cdots\wedge
 c_1(L_n,\|\cdot\|_n).
 \end{aligned}
\]
for model metrics, and hence for continuous semipositive metrics by
uniform approximation.

\section{Dictionary established in this chapter}
\label{sec:metrics-dictionary}
\index{dictionary!metrics}

\[
\begin{array}{ccl}
(\mathcal X,m^{-1}\mathcal L)
 &\longleftrightarrow& \text{model metric on }L,\\[1mm]
\mathcal L_2-\mathcal L_1=D
 &\longleftrightarrow& \text{relative weight }\varphi_D,\\[1mm]
\mathcal L\text{ relatively nef}
 &\longleftrightarrow& \text{semipositive model metric},\\[1mm]
\text{finite weighted sections}
 &\longleftrightarrow& \text{Fubini--Study metric},\\[1mm]
\text{intersection degrees on }\mathcal X_s
 &\longleftrightarrow& \text{atomic Chambert--Loir measure}.
\end{array}
\]

In more detail, and with the place where each entry is proved, the chapter
has established the following correspondences. As in
\cref{sec:models-dictionary}, the table asserts an equality of constructions,
not an analogy.

\medskip
\noindent
\begin{tabular}{@{}p{.33\textwidth}p{.34\textwidth}p{.25\textwidth}@{}}
\toprule
model or section datum & metric or measure datum & where\\
\midrule
line bundle \(\mathcal L\) on a model
 & continuous metric on \(L^{\mathrm{an}}\)
 & \Cref{prop:metrics-model-continuous}\\
vertical Cartier divisor \(D\)
 & relative weight \(\varphi_D=v_x(f_D)\)
 & \Cref{prop:metric-difference-model}\\
special fiber \(\mathcal X_s\)
 & constant weight \(1\); scaling by \(\rme^{-1}\)
 & \Cref{prop:metrics-divisor-model-function}\,(3)\\
\(D\ge0\) effective
 & \(\varphi_D\ge0\), tested at the \(x_E\)
 & \Cref{prop:metrics-divisor-model-function}\,(4)\\
sum of two vertical ideals
 & minimum of two model functions
 & \Cref{lem:metrics-model-lattice}\\
base-point-free sections, weights in \(\Gamma_{\mathbf Q}\)
 & semipositive model metric
 & \Cref{prop:metrics-fs-model}\\
arbitrary real weights
 & continuous semipositive metric
 & \Cref{prop:metrics-fs-real-weights}\\
model line bundles at all levels
 & uniformly dense set of metrics
 & \Cref{prop:metrics-model-function-density}\\
\(\mathcal L\) relatively nef
 & semipositive metric, model independent
 & \Cref{lem:metrics-nef-independence}\\
pullback, tensor, maximum
 & the same operations on metrics
 & \Cref{prop:semipositive-stability}\\
degrees \(b_Y(\mathcal L_1\cdots\mathcal L_n|_Y)\)
 & atoms of the measure at \(x_Y\)
 & \Cref{thm:cl-model-properties}\\
vertical intersection number
 & \(\int\varphi_D\,c_1(L_1)\cdots c_1(L_n)\)
 & \Cref{prop:metrics-cl-model-function-integral}\\
generic degree \((L_1\cdots L_n)\)
 & total mass of the measure
 & \Cref{thm:cl-model-properties}\\
\bottomrule
\end{tabular}
\medskip

Four computations were carried out in full and should be remembered as the
test suite for later chapters: the standard and once-blown-up models of
\(\mathbf P^1\) (\Cref{ex:metrics-p1-standard-model},
\Cref{ex:metrics-p1-blowup}, \Cref{ex:metrics-cl-p1-blowup}); a
Fubini--Study metric with two corners
(\Cref{ex:metrics-fs-two-corners}); a maximum of two boundary-semipositive
metrics on the conic skeleton, computed both by Laplacians and by
intersection numbers (\Cref{ex:metrics-max-skeleton}); and a toric surface,
where the same measure is obtained from degrees on the two components of a
special fiber and from the volume of a subgradient polytope
(\Cref{ex:metrics-cl-toric-surface}). In every one of them the total mass
equals the generic intersection number, as it must.

\Cref{part:cld} will reinterpret the model measure as a wedge product of
Chern currents. \Cref{part:bfj} will extend it from continuous to singular
finite-energy metrics. \Cref{part:bj} will replace the reference
\(v(\pi)=1\) by homogeneity and rebuild the Fubini--Study layer as a
primitive notion. The master comparison is assembled in
\cref{chap:master-dictionary}.

\section*{Source notes}

Model and continuous semipositive metrics originate in Zhang's work
\cite{Zhang1995}.
Chambert--Loir introduced the atomic intersection measure and its continuous
extension in \cite{CL06}. The discrete, residue-characteristic-zero
model-function formalism used by BFJ is in \cite[Sections~2--5]{BFJ16}.
Gubler--Martin \cite{GM19} prove, over arbitrary complete nontrivially valued fields, the
equivalence among algebraic, formal, and PL metrics on proper schemes, the
density of piecewise \(\mathbf Q\)-linear metrics, and the closedness theorem
quoted above. Gubler--Künnemann supply the delta-form positivity comparison;
Gubler--Rabinoff supply the newer curve-test framework. Those comparisons
are stated with their exact hypotheses in
\crefrange{chap:cld-metrized-bundles}{chap:cld-examples} and are not used
here to enlarge the definition silently.

A few words on how the subject reached the form presented here. Metrics
defined by integral models appear in Arakelov theory well before Berkovich
spaces were available: at a finite place of a number field, an integral
model of a projective variety defines a norm on the fiber of a line bundle
at every algebraic point, and heights are computed from such data. Zhang's
paper \cite{Zhang1995} isolated the two features that survive
non-Archimedean analytification: that these metrics can be added, subtracted,
and uniformly approximated, and that relative nefness of the model line
bundle is the right positivity condition, stable enough to define a class of
continuous semipositive metrics by approximation. What was missing was a
space on which such a metric is a genuinely continuous object. Berkovich's
theory supplied it, and Chambert-Loir's article \cite{CL06} then constructed
the measures: the atomic formula of \cref{sec:metrics-cl-measure} is his, as
is the continuous extension of \Cref{thm:cl-continuous-extension}. The
equidistribution applications for which that paper was written are outside
the scope of this book, but the local-height recursion it establishes is
used verbatim in the proof given above.

The density of model functions in \(C^0(X^{\mathrm{an}})\), our
\Cref{prop:metrics-model-function-density}, is due to Gubler, who proved it
in his work on local heights of subvarieties over non-Archimedean fields;
the reader will find it stated for piecewise linear metrics over an
arbitrary complete nontrivially valued field as
\cite[Theorem~2.17]{GM19} and, in the model-function form used by
Boucksom--Favre--Jonsson, in \cite[\S2]{BFJ16}. The proof given here---the
lattice property, separation by Fubini--Study differences, and the lattice
form of the Stone--Weierstrass theorem---is the standard one, arranged so
that no ingredient outside \cref{chap:berkovich,chap:models} is used.

Gubler--Martin's \cite{GM19} is the reference for everything in
\cref{sec:metrics-formal-pl}, and also for the maximum theorem quoted in
\Cref{prop:semipositive-stability}\,(3) and for the closedness theorem
\Cref{thm:model-semipositive-closed}. Its scope is worth emphasizing,
because it is easy to under-state: it removes discreteness of the valuation
and residue-characteristic hypotheses from statements that had previously
been available only in the BFJ setting, and it does so by working with
formal models and piecewise linear metrics throughout. The curve case of the
maximum theorem, proved here as \Cref{lem:metrics-max-graph}, is much older
and is essentially the observation that a maximum of subharmonic functions
is subharmonic.

The toric picture of \Cref{ex:metrics-cl-toric-surface} belongs to a
systematic theory. Burgos Gil--Philippon--Sombra \cite{BGPS14} develop the
arithmetic geometry of toric varieties, including the correspondence between
toric metrics and concave or convex functions, the identification of
Chambert--Loir measures with real Monge--Ampère measures, and the
computation of heights as integrals over polytopes; \cite{CLS11} and
\cite{Fultor} are the standard references for the underlying toric geometry,
and \cite{BGJKK25} develops the invariant tropical theory to which
\cref{chap:model-calculations} refers. Vilsmeier's comparison
\cite{Vil20} treats the corresponding statement on a polystable face with
its calibration factor, and is quoted in this book only in that form.

Three developments after the foundational period are relevant to later
chapters and are mentioned here so that the reader can date the material.
Gubler--Jell--Künnemann--Martin \cite{GJKM19} proved continuity of
plurisubharmonic envelopes using test ideals, which removes a
residue-characteristic hypothesis from the BFJ theory at a decisive point;
Boucksom--Gubler--Martin \cite{BGM21} and
Burgos Gil--Gubler--Jell--Künnemann--Martin \cite{BGGJKM20} developed
non-Archimedean volumes and their differentiability, which is the
quantitative form of the approximation theory begun in this chapter; and
Fang--Gubler--Künnemann \cite{FGK25} extended the solution of the
non-Archimedean Monge--Ampère equation to mixed characteristic. None of
these is used in the present chapter, and all of them presuppose exactly the
model calculus constructed here.

Finally, a remark on what has deliberately not been done. We have not
defined a psh metric, a singular metric, or a metric of finite energy;
we have not defined any local notion of positivity; and we have not asserted
that the semipositive metrics of \Cref{prop:semipositive-stability} form a
closed convex set in any topology finer than the uniform one. Each of these
requires its own chapter, and each of them is where the theories of
\cref{part:cld}, \cref{part:bfj}, and \cref{part:bj} diverge.
